% Chapter 4, Section 2 _Linear Algebra_ Jim Hefferon
%  http://joshua.smcvt.edu/linearalgebra
%  2001-Jun-12
\section{Keserupaan}
\index{keserupaan|(}
%<*SimilarityMotiviation0>
Kita telah mendefinisikan dua matriks \( H \) dan \( \hat{H} \) sebagai
ekuivalen jika terdapat matriks nonsingular \( P \) dan \( Q \) sedemikian
sehingga \( \hat{H}=PHQ \).
Motivasinya adalah diagram berikut, yang menunjukkan bahwa $H$ dan $\hat{H}$
sama-sama merepresentasikan pemetaan $h$, tetapi terhadap pasangan basis yang
berbeda, yaitu $B,D$ dan $\hat{B},\hat{D}$.
\begin{equation*}
  \begin{CD}
    V_{\wrt{B}}                   @>h>H>        W_{\wrt{D}}       \\
    @V{\scriptstyle\identity} VV             @V{\scriptstyle\identity} VV \\
    V_{\wrt{\hat{B}}}             @>h>\hat{H}>  W_{\wrt{\hat{D}}}
  \end{CD}
\end{equation*}
%</SimilarityMotiviation0>

%<*SimilarityMotiviation1>
Sekarang kita meninjau kasus khusus transformasi, ketika kodomain sama dengan
domain, dan menambahkan syarat bahwa basis kodomain sama dengan basis domain.
Jadi, kita meninjau representasi terhadap $B,B$ dan $D,D$.
\index{diagram panah}
\begin{equation*}
  \begin{CD}
    V_{\wrt{B}}                   @>t>T>        V_{\wrt{B}}       \\
    @V{\scriptstyle\identity} VV              @V{\scriptstyle\identity} VV \\
    V_{\wrt{D}}                   @>t>\hat{T}>        V_{\wrt{D}}
  \end{CD}
\end{equation*}
Dalam bentuk matriks,
\(\rep{t}{D,D}
  =\rep{\identity}{B,D}\;\rep{t}{B,B}\;\bigl(\rep{\identity}{B,D}\bigr)^{-1} \).
%</SimilarityMotiviation1>



\subsection{Definisi dan Contoh}
\begin{example}
Tinjau transformasi turunan
$\map{d/dx}{\polyspace_2}{\polyspace_2}$,
dan dua basis bagi ruang tersebut,
$B=\sequence{x^2,x,1}$ dan
$D=\sequence{1,1+x,1+x^2}$.
Kita akan menghitung keempat sisi persegi panahnya.
\begin{equation*}
  \begin{CD}
    {\polyspace_2\,}_{\wrt{B}}                   @>d/dx>T>        {\polyspace_2\,}_{\wrt{B}}       \\
    @V{\scriptstyle\identity} VV              @V{\scriptstyle\identity} VV \\
    {\polyspace_2\,}_{\wrt{D}}                   @>d/dx>\hat{T}>        {\polyspace_2\,}_{\wrt{D}}
  \end{CD}
\end{equation*}
Mulailah dari sisi atas.
Pengaruh transformasi pada basis awal~$B$
\begin{equation*}
  x^2\mapsunder{d/dx} 2x
  \qquad
  x\mapsunder{d/dx} 1
  \qquad
  1\mapsunder{d/dx} 0
\end{equation*}
yang direpresentasikan terhadap basis akhir (juga~$B$)
\begin{equation*}
  \rep{2x}{B}=\colvec{0 \\ 2 \\ 0}
  \qquad
  \rep{1}{B}=\colvec{0 \\ 0 \\ 1}
  \qquad
  \rep{0}{B}=\colvec{0 \\ 0 \\ 0}
\end{equation*}
memberikan representasi pemetaannya.
\begin{equation*}
  T=
  \rep{d/dx}{B,B}=
  \begin{mat}
    0 &0 &0 \\
    2 &0 &0 \\
    0 &1 &0
  \end{mat}
\end{equation*}

Selanjutnya, menghitung matriks bagi sisi kanan melibatkan pencarian pengaruh
pemetaan identitas pada unsur-unsur~$B$.
Tentu saja, pemetaan identitas sama sekali tidak mengubahnya, sehingga untuk
mencari matriksnya kita merepresentasikan unsur-unsur $B$ terhadap~$D$.
\begin{equation*}
  \rep{x^2}{D}=\colvec{-1 \\ 0 \\ 1}
  \quad
  \rep{x}{D}=\colvec{-1 \\ 1 \\ 0}
  \quad
  \rep{1}{D}=\colvec{1 \\ 0 \\ 0}
\end{equation*}
Jadi, matriks untuk bergerak turun pada sisi kanan merupakan penggabungan
vektor-vektor tersebut.
\begin{equation*}
  P=\rep{id}{B,D}=
  \begin{mat}
    -1 &-1  &1 \\
     0 &1   &0 \\
     1 &0   &0   
  \end{mat}
\end{equation*}

Dengan itu, kita memiliki dua pilihan untuk menghitung matriks bagi gerak naik
pada sisi kiri.
Perhitungan langsung merepresentasikan unsur-unsur~$D$ terhadap~$B$
\begin{equation*}
  \rep{1}{B}=\colvec{0 \\ 0 \\ 1}
  \quad
  \rep{1+x}{B}=\colvec{0 \\ 1 \\ 1}
  \quad
  \rep{1+x^2}{B}=\colvec{1 \\ 0 \\ 1}
\end{equation*}
dan menggabungkannya untuk membentuk matriks.
\begin{equation*}
  \begin{mat}
    0 &0 &1 \\
    0 &1 &0 \\ 
    1 &1 &1
  \end{mat}
\end{equation*}
Pilihan lain untuk menghitung matriks bagi gerak naik pada sisi kiri adalah
mengambil invers matriks~$P$ bagi gerak turun pada sisi kanan.
\begin{equation*}
  \begin{pmat}{ccc|ccc}
    -1 &-1 &1 &1 &0 &0 \\
     0 &1  &0 &0 &1 &0 \\
     1 &0  &0 &0 &0 &1
  \end{pmat}
  % \grstep{\rho_1+\rho_3}
  % \grstep{\rho_2+\rho_3}
  % \grstep{-\rho_1}
  % \grstep{\rho_3+\rho_1}
  % \grstep{-\rho_2+\rho_1}
  \grstep{}
  \cdots
  \grstep{}
  \begin{pmat}{ccc|ccc}
     1 &0  &0 &0 &0 &1 \\
     0 &1  &0 &0 &1 &0 \\
     0 &0  &1 &1 &1 &1
  \end{pmat}
\end{equation*}

Tersisa sisi bawah persegi.
Ada dua cara menghitung matriks~$\hat{T}$.
Cara pertama ialah menghitungnya langsung dengan mencari pengaruh transformasi
pada unsur-unsur~$D$
\begin{equation*}
  1\mapsunder{d/dx} 0
  \qquad
  1+x\mapsunder{d/dx} 1
  \qquad
  1+x^2\mapsunder{d/dx} 2x
\end{equation*}
yang direpresentasikan terhadap~$D$.
\begin{equation*}
  \hat{T}=
  \rep{d/dx}{D,D}=
  \begin{mat}
    0 &1 &-2 \\
    0 &0 &2 \\
    0 &0 &0
  \end{mat}
\end{equation*}
Cara lain menghitung~$\hat{T}$, yang biasanya akan kita gunakan, ialah mengikuti
diagram ke atas, menyeberang, lalu turun.
\begin{align*}
  \rep{d/dx}{D,D}
  &=\rep{\identity}{B,D}\,\rep{d/dx}{B,B}\,\rep{\identity}{D,B}  \\
  \hat{T}
  &=\rep{\identity}{B,D}\,T\,\rep{\identity}{D,B}   \\
  &=  
  \begin{mat}
    -1 &-1  &1 \\
     0 &1   &0 \\
     1 &0   &0   
  \end{mat}
  \begin{mat}
    0 &0 &0 \\
    2 &0 &0 \\
    0 &1 &0
  \end{mat}
  \begin{mat}
    0 &0 &1 \\
    0 &1 &0 \\ 
    1 &1 &1
  \end{mat}
\end{align*}
Perkalian tersebut memberikan matriks~$\hat{T}$ yang sama dengan hasil di atas.
\end{example}

\begin{definition} \label{df:Similar}
%<*df:Similar>
Matriks \( T \) dan $\hat{T}$ disebut
\definend{serupa}\index{matriks!keserupaan}%
\index{relasi ekuivalensi!keserupaan matriks}\index{matriks serupa}
jika terdapat matriks nonsingular \( P \) sedemikian sehingga
$
  \hat{T}=PTP^{-1}
$.
%</df:Similar>
\end{definition}

\noindent Karena matriks nonsingular berbentuk persegi, $T$ dan $\hat{T}$ harus
berbentuk persegi dengan ukuran yang sama.
\nearbyexercise{exer:SimIsEquivRel} memeriksa bahwa keserupaan merupakan relasi
ekuivalensi.

\begin{example}
Definisi tersebut tidak mengharuskan kita meninjau suatu pemetaan.
Perhitungan dengan dua matriks berikut,
\begin{equation*}
  P=
  \begin{mat}[r]
    2  &1  \\
    1  &1
  \end{mat}
  \qquad
  T=
  \begin{mat}[r]
    2  &-3  \\
    1  &-1
  \end{mat}
\end{equation*}
menunjukkan bahwa $T$ serupa dengan matriks berikut.
\begin{equation*}
  \hat{T}=
  \begin{mat}[r]
    12  &-19  \\
    7  &-11
  \end{mat}
\end{equation*}
\end{example}

\begin{example}  \label{ex:OnlyZeroSimToZero}
%<*ex:OnlyZeroSimToZero>
Satu-satunya matriks yang serupa dengan matriks nol adalah matriks nol itu
sendiri:~$PZP^{-1}=PZ=Z$.
Matriks identitas memiliki sifat yang sama:~$PIP^{-1}=PP^{-1}=I$.
%</ex:OnlyZeroSimToZero>
\end{example}

Kasus khusus yang umum terjadi ketika ruang vektornya adalah $\C^n$ dan
matriks~$T$ merepresentasikan suatu pemetaan terhadap basis standar.
\begin{equation*}
  \begin{CD}
    \C^n_{\wrt{E_n}}                   @>t>T>        \C^n_{\wrt{E_n}}       \\
    @V{\scriptstyle\identity} VV              @V{\scriptstyle\identity} VV \\
    \C^n_{\wrt{D}}                   @>t>\hat{T}>        \C^n_{\wrt{D}}
  \end{CD}
\end{equation*}
Dalam kasus ini, pada persamaan keserupaan $\hat{T}=PTP^{-1}$, kolom-kolom~$P$
merupakan unsur-unsur~$D$.

Keserupaan matriks merupakan kasus khusus ekuivalensi matriks, sehingga jika
dua matriks serupa, keduanya ekuivalen.
Bagaimana dengan kebalikannya:~jika keduanya berbentuk persegi, apakah sebarang
dua matriks ekuivalen harus serupa?
Tidak; kelas ekuivalensi matriks dari suatu matriks identitas terdiri atas semua
matriks nonsingular berukuran sama, sedangkan contoh sebelumnya menunjukkan
bahwa satu-satunya anggota kelas keserupaan matriks identitas adalah matriks itu
sendiri.
Jadi, kedua matriks berikut ekuivalen tetapi tidak serupa.
\begin{equation*}
  T=
  \begin{mat}[r]
    1  &0  \\
    0  &1
  \end{mat}
  \qquad
  S=
  \begin{mat}[r]
    1  &2  \\
    0  &3
  \end{mat}
\end{equation*}
Jadi, sebagian kelas ekuivalensi matriks terpecah menjadi dua atau lebih kelas
keserupaan\Dash keserupaan memberikan partisi yang lebih halus daripada
ekuivalensi matriks.
Gambar berikut menunjukkan beberapa kelas ekuivalensi matriks yang terbagi
menjadi kelas-kelas keserupaan.
\begin{center}
  \includegraphics{jc/mp/ch5.4}
\end{center}

Untuk memahami relasi keserupaan, kita akan mempelajari kelas-kelas keserupaan.
Kita mendekati pertanyaan ini seperti ketika mempelajari relasi ekuivalensi
baris dan ekuivalensi matriks, yaitu dengan mencari bentuk kanonik bagi
wakil-wakil % \appendrefs{representatives}\spacefactor=1000 %
kelas keserupaan, yang disebut bentuk Jordan.
Dengan bentuk kanonik ini, kita dapat menentukan apakah dua matriks serupa
dengan memeriksa apakah keduanya berada dalam kelas yang memiliki wakil sama.
Dalam ekuivalensi baris maupun ekuivalensi matriks, kita juga telah melihat
bahwa bentuk kanonik memberi wawasan tentang kesamaan anggota-anggota suatu
kelas (misalnya, dua matriks berukuran sama ekuivalen jika dan hanya jika
keduanya memiliki peringkat yang sama).
%(Along the way we shall see ideas that are interesting and important
%in their own right, not just as stepping stones to Jordan form.)

\begin{exercises}
  \item 
    Untuk
    \begin{equation*}
      T=
      \begin{mat}[r]
        1  &3  \\
       -2  &-6
      \end{mat}
      \quad
      \hat{T}=
      \begin{mat}[r]
        0    &0  \\
       -11/2 &-5
      \end{mat}
      \quad
      P=
      \begin{mat}[r]
        4  &2  \\
       -3  &2
      \end{mat}
    \end{equation*}
    periksa bahwa $\hat{T}=PTP^{-1}$.
    \begin{answer}
      Salah satu cara mengerjakannya adalah dari kiri ke kanan.
      \begin{multline*}
        PTP^{-1}=
        \begin{mat}[r]
          4  &2  \\
         -3  &2
        \end{mat}
        \begin{mat}[r]
          1  &3  \\
         -2  &-6
        \end{mat}
        \begin{mat}[r]
          2/14  &-2/14  \\
          3/14  &4/14
        \end{mat}                                 \\
        =
        \begin{mat}[r]
          0  &0  \\
         -7  &-21
        \end{mat}  
        \begin{mat}[r]
          2/14  &-2/14  \\
          3/14  &4/14
        \end{mat}
        =
        \begin{mat}[r]
          0    &0  \\
         -11/2 &-5
        \end{mat}
      \end{multline*}
    \end{answer}
  \item
    \nearbyexample{ex:OnlyZeroSimToZero} menunjukkan bahwa satu-satunya matriks
    yang serupa dengan matriks nol adalah matriks nol itu sendiri, dan satu-
    satunya matriks yang serupa dengan matriks identitas adalah matriks identitas
    itu sendiri.
    \begin{exparts}
      \partsitem Tunjukkan bahwa matriks $\nbyn{1}$ yang satu-satunya entrinya
         adalah $2$ juga hanya serupa dengan dirinya sendiri.
      \partsitem Apakah matriks berbentuk $cI$ bagi suatu skalar $c$ hanya serupa
         dengan dirinya sendiri?
     \partsitem Apakah matriks diagonal hanya serupa dengan dirinya sendiri?
     % \partsitem Is a square block partial-identity matrix similar only to 
     % itself?
    \end{exparts}
    \begin{answer}
     \begin{exparts}
      \partsitem Karena matriks $(2)$ berukuran $\nbyn{1}$, matriks $P$ dan
         $P^{-1}$ juga berukuran $\nbyn{1}$. Jika $P=(p)$, inversnya adalah
         $P^{-1}=(1/p)$. Jadi $P(2)P^{-1}=(p)(2)(1/p)=(2)$.
      \partsitem Ya:~ingat bahwa kelipatan skalar dapat dikeluarkan dari suatu
        matriks, \( P(cI)P^{-1}=cPIP^{-1}=cI \).
        Matriks nol dan identitas masing-masing merupakan kasus khusus $c=0$
        dan $c=1$.
      \partsitem Tidak, sebagaimana ditunjukkan contoh berikut.
        \begin{equation*}
           \begin{mat}[r]
              1  &-2  \\
             -1  &1
            \end{mat}
           \begin{mat}[r]
             -1  &0   \\
              0  &-3
           \end{mat}
           \begin{mat}[r]
              -1  &-2   \\
              -1  &-1
           \end{mat}
           =
           \begin{mat}[r]
              -5  &-4   \\
              2   &1
           \end{mat}
        \end{equation*} 
    \end{exparts}  
   \end{answer}
  \recommended \item Tinjau transformasi pada~$\C^3$ berikut,
    \begin{equation*}
      t(\colvec{x \\ y \\ z})=\colvec{x-z \\ z \\ 2y}
    \end{equation*}
    beserta basis-basis berikut.
    \begin{equation*}
      B=\sequence{\colvec{1 \\ 2 \\ 3}, 
                  \colvec{0 \\ 1 \\ 0}, 
                  \colvec{0 \\ 0 \\ 1}}
      \qquad
      D=\sequence{\colvec{1 \\ 0 \\ 0},
                  \colvec{1 \\ 1 \\ 0},
                  \colvec{1 \\ 0 \\ 1}}
    \end{equation*}
    Kita akan menghitung bagian-bagian diagram panah untuk merepresentasikan
    transformasi tersebut menggunakan dua matriks serupa.
    \begin{exparts}
      \partsitem Gambarkan diagram panah yang dikhususkan bagi kasus ini.
      \partsitem Hitung $T=\rep{t}{B,B}$.
      \partsitem Hitung $\hat{T}=\rep{t}{D,D}$.
      \partsitem Hitung matriks bagi dua sisi lain persegi panah.
    \end{exparts}
    \begin{answer}
      \begin{exparts}
        \partsitem
          \begin{equation*}
            \begin{CD}
              \C^3_{\wrt{B}}                   @>t>T>        \C^3_{\wrt{B}}       \\
              @V{\scriptstyle\identity} VV              @V{\scriptstyle\identity} VV \\
              \C^3_{\wrt{D}}                   @>t>\hat{T}>        \C^3_{\wrt{D}}
            \end{CD}
          \end{equation*}
        \partsitem
          Bagi setiap unsur basis awal~$B$, cari pengaruh transformasinya,
          \begin{equation*}
             \colvec{1 \\ 2 \\ 3}\mapsunder{t}\colvec{-2 \\ 3 \\ 4}
             \qquad 
             \colvec{0 \\ 1 \\ 0}\mapsunder{t}\colvec{0 \\ 0 \\ 2}
             \qquad 
             \colvec{0 \\ 0 \\ 1}\mapsunder{t}\colvec{-1 \\ 1 \\ 0}
          \end{equation*}
          lalu representasikan hasil-hasilnya terhadap basis akhir~$B$,
          \begin{equation*}
            \rep{\colvec{-2 \\ 3 \\ 4}}{B}=\colvec{-2 \\ 7 \\ 10}
            \qquad
            \rep{\colvec{0 \\ 0 \\ 2}}{B}=\colvec{0 \\ 0 \\ 2}
            \qquad
            \rep{\colvec{-1 \\ 1 \\ 0}}{B}=\colvec{-1 \\ 3 \\ 3}
          \end{equation*}
          untuk memperoleh matriks.
          \begin{equation*}
            T=\rep{t}{B,B}=
            \begin{mat}
              -2 &0 &-1 \\
               7 &0 &3  \\
              10 &2 &3 
            \end{mat}
          \end{equation*}
        \partsitem
          Cari pengaruh transformasi pada unsur-unsur~$D$,
          \begin{equation*}
             \colvec{1 \\ 0 \\ 0}\mapsunder{t}\colvec{1 \\ 0 \\ 0}
             \qquad 
             \colvec{1 \\ 1 \\ 0}\mapsunder{t}\colvec{1 \\ 0 \\ 2}
             \qquad 
             \colvec{1 \\ 0 \\ 1}\mapsunder{t}\colvec{0 \\ 1 \\ 0}
          \end{equation*}
          lalu representasikan hasil-hasilnya terhadap basis akhir~$D$,
          \begin{equation*}
            \rep{\colvec{1 \\ 0 \\ 0}}{D}=\colvec{1 \\ 0 \\ 0}
            \qquad
            \rep{\colvec{1 \\ 0 \\ 2}}{D}=\colvec{-1 \\ 0 \\ 2}
            \qquad
            \rep{\colvec{0 \\ 1 \\ 0}}{D}=\colvec{-1 \\ 1 \\ 0}
          \end{equation*}
          untuk memperoleh matriks.
          \begin{equation*}
            \hat{T}=\rep{t}{D,D}=
            \begin{mat}
               1 &-1 &-1 \\
               0 &0  &1  \\
               0 &2  &0 
            \end{mat}
          \end{equation*}
        \partsitem
          Untuk bergerak turun pada sisi kanan, kita memerlukan
          $\rep{\identity}{B,D}$. Jadi, mula-mula hitung pengaruh pemetaan
          identitas pada setiap unsur~$B$, yang tidak mengubah unsur tersebut,
          lalu representasikan hasilnya terhadap~$D$.
          \begin{equation*}
            \rep{\colvec{1 \\ 2 \\ 3}}{D}=\colvec{-4 \\ 2 \\ 3}
            \qquad
            \rep{\colvec{0 \\ 1 \\ 0}}{D}=\colvec{-1 \\ 1 \\ 0}
            \qquad
            \rep{\colvec{0 \\ 0 \\ 1}}{D}=\colvec{-1 \\ 0 \\ 1}
          \end{equation*}
          Jadi, inilah~$P$.
          \begin{equation*}
            P=
            \begin{mat}
              -4 &-1 &-1 \\
               2 &1  &0  \\
               3 &0  &1
            \end{mat}
          \end{equation*}
          Matriks lain~$\rep{\identity}{D,B}$ dapat dicari secara langsung
          seperti ketika mencari~$P$, atau dengan menghitung invers matriks
          seperti biasa.
          \begin{equation*}
            P^{-1}=
            \begin{mat}
               1 &1 &1 \\
               -2 &-1  &-2  \\
               -3 &-3  &-2
            \end{mat}
          \end{equation*}
      \end{exparts}
    \end{answer}
  \item 
    Tinjau transformasi $\map{t}{\polyspace_2}{\polyspace_2}$ yang
    dideskripsikan oleh $x^2\mapsto x+1$, $x\mapsto x^2-1$, dan $1\mapsto 3$.
    \begin{exparts}
      \partsitem Cari $T=\rep{t}{B,B}$ dengan $B=\sequence{x^2,x,1}$.
      \partsitem Cari $\hat{T}=\rep{t}{D,D}$ dengan
        $D=\sequence{1,1+x,1+x+x^2}$.
      \partsitem Cari matriks $P$ sedemikian sehingga $\hat{T}=PTP^{-1}$.
    \end{exparts}
    \begin{answer}
      \begin{exparts}
        \partsitem Karena kita mendeskripsikan $t$ menggunakan anggota-anggota
          $B$, mencari representasi matriksnya mudah:
          \begin{equation*}
            \rep{t(x^2)}{B}=\colvec[r]{0 \\ 1 \\ 1}_B
            \quad
            \rep{t(x)}{B}=\colvec[r]{1 \\ 0 \\ -1}_B
            \quad
            \rep{t(1)}{B}=\colvec[r]{0 \\ 0 \\ 3}_B
          \end{equation*}
          memberikan hasil berikut.
          \begin{equation*}
            \rep{t}{B,B}=
            \begin{mat}[r]
              0  &1  &0  \\
              1  &0  &0  \\
              1  &-1 &3  
            \end{mat}
          \end{equation*}
        \partsitem Kita akan mencari $t(1)$, $t(1+x)$, dan $t(1+x+x^2)$ untuk
          mengetahui representasi masing-masing terhadap $D$.
          Diberikan $t(1)=3$, dan dua lainnya mudah diperoleh:
          $t(1+x)=x^2+2$ dan $t(1+x+x^2)=x^2+x+3$.
          Dengan mengamati langsung, kita memperoleh representasi setiap vektor,
          \begin{equation*}
            \rep{t(1)}{D}=\colvec[r]{3 \\ 0 \\ 0}_D
            \quad
            \rep{t(1+x)}{D}=\colvec[r]{2  \\ -1 \\  1}_D
            \quad
            \rep{t(1+x+x^2)}{D}=\colvec[r]{2 \\ 0 \\ 1}_D
          \end{equation*}
          sehingga diperoleh representasi pemetaannya.
          \begin{equation*}
            \rep{t}{D,D}
            =
            \begin{mat}[r]
              3  &2  &2  \\
              0  &-1 &0  \\
              0  &1  &1
            \end{mat}
          \end{equation*}
         \partsitem Diagram
           \begin{equation*}
             \begin{CD}
               V_{\wrt{B}}                  @>t>T>  V_{\wrt{B}}       \\
               @V\scriptstyle\identity VPV      @V\scriptstyle\identity VPV \\
               V_{\wrt{D}}                  @>t>\hat{T}>  V_{\wrt{D}}
             \end{CD}
           \end{equation*}
           menunjukkan bahwa matriks-matriks ini adalah
           $P=\rep{\identity}{B,D}$ dan $P^{-1}=\rep{\identity}{D,B}$.
           \begin{equation*}
             P=
             \begin{mat}[r]
               0  &-1  &1  \\ 
               -1  &1  &0  \\
               1  &0  &0 
             \end{mat}
             \qquad
             P^{-1}=
             \begin{mat}[r]
               0  &0  &1  \\ 
               0  &1  &1  \\
               1  &1  &1 
             \end{mat}
           \end{equation*}
      \end{exparts}
    \end{answer}
  \recommended \item 
    Misalkan $T$ merepresentasikan $\map{t}{\C^2}{\C^2}$ terhadap $B,B$.
    \begin{equation*}
      T=
      \begin{mat}
        1 &-1 \\ 
        2 &1
      \end{mat}
      \qquad
      B=\sequence{\colvec{1 \\ 0},\colvec{1 \\ 1}},\hspace{0.7em}
      D=\sequence{\colvec{2 \\ 0},\colvec{0 \\ -2}}
    \end{equation*}
    Kita akan mengubahnya menjadi matriks yang merepresentasikan~$t$ terhadap
    $D,D$.
    \begin{exparts}
      \partsitem Gambarkan diagram panahnya.
      \partsitem Berikan matriks yang merepresentasikan sisi kiri dan kanan
          diagram tersebut, dalam arah yang kita lalui untuk melakukan
          perubahan representasi.
      \partsitem Cari $\rep{t}{D,D}$.
    \end{exparts}
    \begin{answer}
      \begin{exparts}
        \partsitem
          \begin{equation*}
            \begin{CD}
              \C^2_{\wrt{B}}                   @>t>T>        \C^2_{\wrt{B}}       \\
              @V{\scriptstyle\identity} VV              @V{\scriptstyle\identity} VV \\
              \C^2_{\wrt{D}}                   @>t>\hat{T}>        \C^2_{\wrt{D}}
            \end{CD}
          \end{equation*}  
        \partsitem
          Bagi sisi kanan, cari pengaruh pemetaan identitas, yang tidak mengubah
          unsur,
          \begin{equation*}
            \colvec{1 \\ 0}\mapsunder{\identity}\colvec{1 \\ 0}
            \qquad
            \colvec{1 \\ 1}\mapsunder{\identity}\colvec{1 \\ 1}
          \end{equation*}
          lalu representasikan hasilnya terhadap~$D$,
          \begin{equation*}
            \rep{\colvec{1 \\ 0}}{D}=\colvec{1/2 \\ 0}
            \qquad
            \rep{\colvec{1 \\ 1}}{D}=\colvec{1/2 \\ -1/2}
          \end{equation*}
          sehingga diperoleh
          \begin{equation*}
            P=\rep{\identity}{B,D}=
            \begin{mat}
              1/2 &1/2 \\
              0   &-1/2
            \end{mat}
          \end{equation*}

          Matriks pada sisi kiri dapat dihitung secara langsung seperti pada
          paragraf sebelumnya, atau diperoleh dengan mengambil invers.
          \begin{equation*}
            P^{-1}=\rep{\identity}{D,B}=
            \frac{1}{(-1/4)}\cdot
            \begin{mat}
              -1/2 &-1/2 \\
              0   &1/2
            \end{mat}
            =
            \begin{mat}
              2 &2 \\
              0 &-2
            \end{mat}
          \end{equation*}
        \partsitem
          Seperti pada butir sebelumnya, kita dapat menghitungnya langsung dari
          definisi atau menggunakan operasi matriks.
          \begin{equation*}
            PTP^{-1}=
            \begin{mat}
              2 &2 \\
              0 &-2
            \end{mat}
            \begin{mat}
              1 &-1 \\ 
              2 &1
            \end{mat}
            \begin{mat}
              2 &2 \\
              0 &-2
            \end{mat}
            =
            \begin{mat}
              3 &3  \\
             -2 &-1
            \end{mat}
          \end{equation*}
      \end{exparts}
    \end{answer}
  \recommended \item
     Tampilkan suatu hubungan keserupaan tak trivial dengan membiarkan
     \( \map{t}{\C^2}{\C^2} \) bertindak sebagai berikut,
     \begin{equation*}
        \colvec[r]{1 \\ 2}\mapsto\colvec[r]{3 \\ 0}
        \qquad
        \colvec[r]{-1 \\ 1}\mapsto\colvec[r]{-1 \\ 2}
     \end{equation*}
     memilih dua basis~$B,D$, lalu merepresentasikan \( t \) terhadap keduanya,
     \( \hat{T}=\rep{t}{B,B} \) dan \( T=\rep{t}{D,D} \).
     Kemudian hitung \( P \) dan \( P^{-1} \) untuk mengubah basis dari \( B \)
     ke \( D \), lalu kembali lagi.
     \begin{answer}
       Salah satu pilihan basis yang mungkin adalah
       \begin{equation*}
         B=\sequence{\colvec[r]{1 \\ 2},\colvec[r]{-1 \\ 1}}
         \qquad
         D=\stdbasis_2=\sequence{\colvec[r]{1 \\ 0},\colvec[r]{0 \\ 1}}
       \end{equation*}
       (basis $B$ ini berasal dari deskripsi pemetaan).
       Untuk mencari matriks $\hat{T}=\rep{t}{B,B}$, selesaikan hubungan
       \begin{equation*}
          c_1\colvec[r]{1 \\ 2}+c_2\colvec[r]{-1 \\ 1}=\colvec[r]{3 \\ 0}
          \qquad
         \hat{c}_1\colvec[r]{1 \\ 2}+\hat{c}_2\colvec[r]{-1 \\ 1}=\colvec[r]{-1 \\ 2}
       \end{equation*}
       untuk memperoleh \( c_1=1 \), \( c_2=-2 \), \( \hat{c}_1=1/3 \), dan
       \( \hat{c}_2=4/3 \).
       \begin{equation*}
          \rep{t}{B,B}=
          \begin{mat}[r]
             1  &1/3 \\
            -2  &4/3
          \end{mat}
       \end{equation*}

       Mencari \( \rep{t}{D,D} \) memerlukan sedikit lebih banyak perhitungan.
       Mula-mula cari \( t(\vec{e}_1) \).
       Hubungan
       \begin{equation*}
         c_1\colvec[r]{1 \\ 2}+c_2\colvec[r]{-1 \\ 1}=\colvec[r]{1 \\ 0}
       \end{equation*}
       memberikan \( c_1=1/3 \) dan \( c_2=-2/3 \), sehingga
       \begin{equation*}
         \rep{\vec{e}_1}{B}=\colvec[r]{1/3 \\ -2/3}_B
       \end{equation*}
       dan karena itu
       \begin{equation*}
          \rep{t(\vec{e}_1)}{B}=
                 \begin{mat}[r]
                    1  &1/3  \\
                   -2  &4/3
                 \end{mat}_{B,B}
                 \colvec[r]{1/3 \\ -2/3}_B
                 =
                 \colvec[r]{1/9 \\ -14/9}_B
       \end{equation*}
       Jadi, $t$ bertindak pada vektor basis pertama $\vec{e}_1$ sebagai
       berikut.
       \begin{equation*}
         t(\vec{e}_1)
         =(1/9)\cdot\colvec[r]{1 \\ 2}-(14/9)\cdot\colvec[r]{-1 \\ 1}
         =\colvec[r]{5/3 \\ -4/3}      
       \end{equation*}
       Perhitungan bagi \( t(\vec{e}_2) \) serupa.
       Hubungan
       \begin{equation*}
         c_1\colvec[r]{1 \\ 2}+c_2\colvec[r]{-1 \\ 1}=\colvec[r]{0 \\ 1}
       \end{equation*}
       memberikan \( c_1=1/3 \) dan \( c_2=1/3 \), sehingga
       \begin{equation*}
         \rep{\vec{e}_2}{B}=\colvec[r]{1/3 \\ 1/3}_B
       \end{equation*}
       dan karena itu
       \begin{equation*}
         \rep{t(\vec{e}_2)}{B}=
                 \begin{mat}[r]
                    1  &1/3  \\
                   -2  &4/3
                 \end{mat}_{B,B}
                 \colvec[r]{1/3 \\ 1/3}_B
                 =
                 \colvec[r]{4/9 \\ -2/9}_B
       \end{equation*}
       Jadi, $t$ bertindak pada vektor basis kedua $\vec{e}_2$ sebagai berikut.
       \begin{equation*}
         t(\vec{e}_2)
         =(4/9)\cdot\colvec[r]{1 \\ 2}-(2/9)\cdot\colvec[r]{-1 \\ 1}
         =\colvec[r]{2/3 \\ 2/3}
       \end{equation*}
       Karena itu,
       \begin{equation*}
          \rep{t}{D,D}=
          \begin{mat}[r]
             5/3  &2/3  \\
            -4/3  &2/3
          \end{mat}
       \end{equation*}
       Matriks perubahan basisnya adalah
       \begin{equation*}
         P=\rep{\identity}{B,D}
         =\begin{mat}[r]
           1  &-1  \\
           2  &1 
         \end{mat}
       \end{equation*}
      sedangkan matriks berikut mengubah basis kembali.
      \begin{equation*}
         P^{-1}=\bigl(\rep{\identity}{B,D}\bigr)^{-1}
         =\begin{mat}[r]
            1  &-1 \\
            2  &1
         \end{mat}^{-1}
         =
         \begin{mat}[r]
           1/3  &1/3  \\
           -2/3 &1/3
         \end{mat}
      \end{equation*}
       Pemeriksaan perhitungan ini bersifat rutin.
      \begin{equation*}
         \begin{mat}[r]
            1  &-1 \\
            2  &1
         \end{mat}
         \begin{mat}[r]
            1  &1/3 \\
           -2  &4/3
         \end{mat}
         \begin{mat}[r]
           1/3 &1/3 \\
          -2/3 &1/3
         \end{mat}
         =
         \begin{mat}[r]
           5/3 &2/3 \\
          -4/3 &2/3
         \end{mat}
      \end{equation*}  
    \end{answer}
  \recommended \item 
    Tunjukkan bahwa matriks-matriks berikut tidak serupa.
    \begin{equation*}
       \begin{mat}[r]
          1  &0  &4  \\
          1  &1  &3  \\
          2  &1  &7
       \end{mat}
       \qquad
       \begin{mat}[r]
          1  &0  &1  \\
          0  &1  &1  \\
          3  &1  &2
       \end{mat}
    \end{equation*}
    \begin{answer}
      Metode Gauss menunjukkan bahwa matriks pertama merepresentasikan pemetaan
      berperingkat dua, sedangkan matriks kedua merepresentasikan pemetaan
      berperingkat tiga.
    \end{answer}
  \item 
    Jelaskan \nearbyexample{ex:OnlyZeroSimToZero} dalam kerangka pemetaan.
    \begin{answer}
      Satu-satunya representasi pemetaan nol adalah matriks nol, apa pun
      pasangan basisnya, $\rep{z}{B,D}=Z$. Jadi, khususnya, bagi sebarang basis
      tunggal $B$ kita memiliki $\rep{z}{B,B}=Z$.
      Kasus identitas sedikit berbeda:~satu-satunya representasi pemetaan
      identitas terhadap sebarang $B,B$ adalah identitas
      $\rep{\identity}{B,B}=I$.
      (\textit{Catatan:}~tentu saja, kita telah melihat contoh dengan $B\neq D$
      dan $\rep{\identity}{B,D}\neq I$\Dash bahkan, kita telah melihat bahwa
      sebarang matriks nonsingular merupakan representasi pemetaan identitas
      terhadap suatu $B,D$.)
    \end{answer}
  \recommended \item 
    \cite{Halmos}
    Adakah dua matriks \( A \) dan \( B \) yang serupa, sedangkan \( A^2 \) dan
    \( B^2 \) tidak serupa?
    \begin{answer}
      Tidak.
      Jika \( A=PBP^{-1} \), maka
      \( A^2=(PBP^{-1})(PBP^{-1})=PB^2P^{-1} \).
    \end{answer}
  \recommended \item
    Buktikan bahwa jika dua matriks serupa dan salah satunya dapat diinverskan,
    matriks lainnya juga dapat diinverskan.
    \begin{answer}
       Keserupaan matriks merupakan kasus khusus ekuivalensi matriks (jika dua
       matriks serupa, keduanya ekuivalen), dan ekuivalensi matriks
       mempertahankan nonsingularitas.
    \end{answer}
  \item \label{exer:SimIsEquivRel}
    Tunjukkan bahwa keserupaan merupakan relasi ekuivalensi.
    (Definisi yang diberikan sebelumnya sudah mencerminkan hal ini. Karena itu,
    mulailah di sini dengan definisi bahwa $\hat{T}$ serupa dengan $T$ jika
    $\hat{T}=PTP^{-1}$.)
    \begin{answer}
       Suatu matriks serupa dengan dirinya sendiri; ambil \( P \) sebagai
       matriks identitas:~$T=ITI^{-1}=ITI$.

       Jika \( \hat{T} \) serupa dengan \( T \), maka
       \( \hat{T}=PTP^{-1} \), sehingga \( P^{-1}\hat{T}P=T \).
       Tulis ulang sebagai \( T=(P^{-1})\hat{T}(P^{-1})^{-1} \) untuk
       menyimpulkan bahwa $T$ serupa dengan \( \hat{T} \).

       Bagi transitivitas, jika \( T \) serupa dengan \( S \) dan \( S \)
       serupa dengan \( U \), maka \( T=PSP^{-1} \) dan \( S=QUQ^{-1} \).
       Jadi, \( T=PQUQ^{-1}P^{-1}=(PQ)U(PQ)^{-1} \), yang menunjukkan bahwa
       \( T \) serupa dengan \( U \).
     \end{answer}
  \item 
     Tinjau suatu matriks yang merepresentasikan pencerminan terhadap sumbu
     \( x \) dalam \( \Re^2 \), terhadap suatu $B,B$.
     Tinjau pula suatu matriks yang merepresentasikan pencerminan terhadap sumbu
     \( y \), terhadap suatu $D,D$.
     Apakah keduanya harus serupa?
     \begin{answer}
        Misalkan $f_x$ dan $f_y$ merupakan pemetaan pencerminan.
        Bagi sebarang basis \( B \) dan \( D \), matriks
        \( \rep{f_x}{B,B} \) dan \( \rep{f_y}{D,D} \) serupa.
        Pertama, perhatikan bahwa
        \begin{equation*}
           S=\rep{f_x}{\stdbasis_2,\stdbasis_2}=
           \begin{mat}[r]
              1  &0  \\
              0  &-1
           \end{mat}
           \qquad
           T=\rep{f_y}{\stdbasis_2,\stdbasis_2}=
           \begin{mat}[r]
             -1  &0  \\
              0  &1
           \end{mat}
        \end{equation*}
        serupa karena matriks kedua merupakan representasi $f_x$ terhadap basis
        \( A=\sequence{\vec{e}_2,\vec{e}_1} \):
        \begin{equation*}
           \begin{mat}[r]
              1  &0  \\
              0  &-1
           \end{mat}
           =    
           P
           \begin{mat}[r]
             -1  &0  \\
              0  &1
           \end{mat}
           P^{-1}
        \end{equation*}
        dengan $P=\rep{\identity}{A,\stdbasis_2}$.
        \begin{equation*}
          \begin{CD}
            \C^2_{\wrt{A}}                   
               @>f_x>T>        
               V\C^2_{\wrt{A}}       \\
            @V\scriptstyle\identity VPV                
               @V\scriptstyle\identity VPV \\
            \C^2_{\wrt{\stdbasis_2}}         
               @>f_x>S>        
               \C^2_{\wrt{\stdbasis_2}}
          \end{CD}
        \end{equation*}
        Sekarang kesimpulannya mengikuti bagian transitivitas pada
        \nearbyexercise{exer:SimIsEquivRel}.
        
        Kita juga dapat menyelesaikannya tanpa bergantung pada latihan tersebut.
        Tulis
        $\rep{f_x}{B,B}=QSQ^{-1}
          =Q\rep{f_x}{\stdbasis_2,\stdbasis_2}Q^{-1}$
        dan
        $\rep{f_y}{D,D}=RSR^{-1}=R\rep{f_y}{\stdbasis_2,\stdbasis_2}R^{-1}$.
        Menurut persamaan pada paragraf pertama, persamaan pertama menjadi
        $\rep{f_x}{B,B}=QP\rep{f_y}{\stdbasis_2,\stdbasis_2}P^{-1}Q^{-1}$.
        Dengan menulis ulang persamaan kedua sebagai
        $R^{-1}\cdot\rep{f_y}{D,D}\cdot R=\rep{f_y}{\stdbasis_2,\stdbasis_2}$
        dan melakukan substitusi, diperoleh hubungan yang diinginkan,
        \begin{multline*}
          \rep{f_x}{B,B}
          =QP\rep{f_y}{\stdbasis_2,\stdbasis_2}P^{-1}Q^{-1}  \\
          =QPR^{-1}\cdot\rep{f_y}{D,D}\cdot RP^{-1}Q^{-1}
          =(QPR^{-1})\cdot\rep{f_y}{D,D}\cdot (QPR^{-1})^{-1}
        \end{multline*}
        Jadi, matriks \( \rep{f_x}{B,B} \) dan \( \rep{f_y}{D,D} \) serupa.
      \end{answer}
  \item 
     Buktikan bahwa keserupaan matriks mempertahankan peringkat dan determinan.
     Apakah kebalikannya berlaku?
     \begin{answer}
       Kita harus menunjukkan bahwa jika dua matriks serupa, keduanya memiliki
       peringkat dan determinan yang sama.

       Pertama, kita telah menunjukkan bahwa peringkat suatu pemetaan sama
       dengan peringkat sebarang matriks yang merepresentasikannya. Jadi,
       peringkat harus sama bagi semua representasi.
       (Dengan kata lain, peringkat merupakan sifat pemetaan karena peringkat
       adalah dimensi daerah hasilnya.)

       Bagi determinan,
       \( \deter{PSP^{-1}}=\deter{P}\cdot\deter{S}\cdot\deter{P^{-1}}
          =\deter{P}\cdot\deter{S}\cdot\deter{P}^{-1}=\deter{S} \).

       Kebalikan pernyataan tersebut tidak berlaku.
       Ada matriks-matriks dengan determinan sama yang tidak serupa.
       Sebagai contoh, tinjau suatu matriks tak nol dengan determinan nol.
       Matriks itu tidak serupa dengan matriks nol, karena matriks nol hanya
       serupa dengan dirinya sendiri, tetapi determinan keduanya sama.
       Peringkat yang sama juga tidak cukup:~dua matriks diagonal $\nbyn{2}$
       dengan diagonal masing-masing $(1,0)$ dan $(2,0)$ sama-sama berperingkat
       satu tetapi tidak serupa, karena jejaknya berbeda.
     \end{answer}
  \item 
    Adakah kelas ekuivalensi matriks yang hanya memuat satu kelas keserupaan
    matriks?
    Adakah yang memuat tak hingga banyak kelas keserupaan?
    \begin{answer}
      Kelas ekuivalensi matriks yang memuat semua matriks $\nbyn{n}$
      berperingkat nol hanya memuat satu matriks, yaitu matriks nol.
      Karena itu, kelas tersebut hanya memuat satu kelas keserupaan.

      Sebaliknya, kelas ekuivalensi matriks $\nbyn{1}$ berperingkat satu terdiri
      atas matriks-matriks $\nbyn{1}$ \( (k) \) dengan \( k\neq 0 \).
      Bagi sebarang basis \( B \), representasi perkalian dengan skalar \( k \)
      adalah \( \rep{t_k}{B,B}=(k) \), sehingga setiap matriks tersebut menjadi
      satu-satunya anggota kelas keserupaannya.
      Jadi, dalam kasus ini satu kelas ekuivalensi matriks terpecah menjadi tak
      hingga banyak kelas keserupaan.
     \end{answer}
  \item 
    Dapatkah dua matriks diagonal yang berbeda berada dalam kelas keserupaan
    yang sama?
    \begin{answer}
      Ya, matriks-matriks berikut serupa,
      \begin{equation*}
         \begin{mat}[r]
           1  &0  \\
           0  &3
         \end{mat}
         \qquad
         \begin{mat}[r]
           3  &0  \\
           0  &1
         \end{mat}
      \end{equation*}
      sebab jika matriks pertama adalah $\rep{t}{B,B}$ bagi
      $B=\sequence{\vec{\beta}_1,\vec{\beta}_2}$, matriks kedua adalah
      $\rep{t}{D,D}$ bagi
      $D=\sequence{\vec{\beta}_2,\vec{\beta}_1}$.
     \end{answer}
  \recommended \item
    Buktikan bahwa jika dua matriks serupa, pangkat ke-\( k \) keduanya serupa
    ketika \( k>0 \).
    Bagaimana jika \( k\leq 0 \)?
    \begin{answer}
      Pangkat ke-\( k \) keduanya serupa karena, jika setiap matriks
      merepresentasikan pemetaan $t$, pangkat ke-\( k \) merepresentasikan
      \( t^k \), yaitu komposisi $k$ buah $t$.
      (Sebagai contoh, jika $T=\rep{t}{B,B}$, maka
      $T^2=\rep{\composed{t}{t}}{B,B}$.)

      Dalam bentuk komputasional, jika \( T=PSP^{-1} \), maka
      \( T^2=(PSP^{-1})(PSP^{-1})=PS^2P^{-1} \).
      Induksi memperluas hasil ini ke semua pangkat positif.

      Untuk $k=0$, kedua pangkat bernilai matriks identitas dan karena itu
      serupa. Untuk $k<0$, pangkat tersebut hanya terdefinisi jika matriksnya
      dapat diinverskan. Andaikan \( S \) dapat diinverskan dan
      \( T=PSP^{-1} \). Perhatikan bahwa \( T \) juga dapat diinverskan:
      \( T^{-1}=(PSP^{-1})^{-1}=PS^{-1}P^{-1} \),
      dan persamaan yang sama menunjukkan bahwa \( T^{-1} \) serupa dengan
      \( S^{-1} \).
      Pangkat negatif lainnya sekarang mengikuti argumen pada paragraf pertama.
    \end{answer}
  \recommended \item
    Misalkan \( p(x) \) adalah polinomial \( c_nx^n+\cdots+c_1x+c_0 \).
    Tunjukkan bahwa jika \( T \) serupa dengan \( S \), maka
    \( p(T)=c_nT^n+\cdots+c_1T+c_0I \) serupa dengan
    \( p(S)=c_nS^n+\cdots+c_1S+c_0I \).
    \begin{answer}
       Secara konseptual, keduanya merepresentasikan \( p(t) \) bagi suatu
       transformasi \( t \).
       Secara komputasional, kita memiliki
       \begin{align*}
          p(T)
          &=c_n(PSP^{-1})^n+\dots+c_1(PSP^{-1})+c_0I   \\
          &=c_nPS^nP^{-1}+\dots+c_1PSP^{-1}+c_0I   \\
           &=Pc_nS^nP^{-1}+\dots+Pc_1SP^{-1}+Pc_0IP^{-1}   \\
           &=P(c_nS^n+\dots+c_1S+c_0I)P^{-1}
       \end{align*}  
     \end{answer}
  \item 
    Daftarkan semua kelas ekuivalensi matriks bagi matriks \( \nbyn{1} \).
    Daftarkan pula kelas keserupaannya, dan deskripsikan kelas keserupaan mana
    yang termuat dalam setiap kelas ekuivalensi matriks.
    \begin{answer}
      Ada dua kelas ekuivalensi:~(i) kelas matriks berperingkat nol, yang hanya
      memiliki satu anggota,
      $\mathscr{C}_1=\set{(0)}$,
      dan (ii) kelas matriks berperingkat satu, yang memiliki tak hingga banyak
      anggota,
      \( \mathscr{C}_2=\set{(k)\suchthat k\neq0} \).

      Setiap matriks \( \nbyn{1} \) merupakan satu-satunya anggota kelas
      keserupaannya.
      Hal itu karena sebarang transformasi pada ruang berdimensi satu merupakan
      perkalian dengan skalar, yaitu $\map{t_k}{V}{V}$ yang diberikan oleh
      $\vec{v}\mapsto k\cdot\vec{v}$.
      Jadi, bagi sebarang basis \( B=\sequence{\vec{\beta}} \), matriks yang
      merepresentasikan transformasi \( t_k \) terhadap \( B,B \) adalah
      \( (\rep{t_k(\vec{\beta})}{B})=(k) \). 
      
      Jadi, kelas ekuivalensi matriks $\mathscr{C}_1$ memuat satu kelas
      keserupaan, yang terdiri atas matriks $(0)$.
      Kelas ekuivalensi matriks $\mathscr{C}_2$ memuat tak hingga banyak kelas
      keserupaan beranggota tunggal, masing-masing terdiri atas $(k)$ bagi
      $k\neq0$.
    \end{answer}
  \item 
    Apakah keserupaan mempertahankan penjumlahan?
    \begin{answer}
      Tidak.
      Berikut contoh dengan dua pasangan, yang masing-masing terdiri atas dua
      matriks serupa:
      \begin{equation*}
         \begin{mat}[r]
            1  &-1  \\
            1  &2
         \end{mat}
         \begin{mat}[r]
            1  &0   \\
            0  &3
         \end{mat}
         \begin{mat}[r]
            2/3  &1/3   \\
           -1/3  &1/3
         \end{mat}
         =
         \begin{mat}[r]
            5/3  &-2/3  \\
           -4/3  &7/3
         \end{mat}
      \end{equation*}
      dan
      \begin{equation*}
         \begin{mat}[r]
            1  &-2  \\
           -1  &1
         \end{mat}
         \begin{mat}[r]
           -1  &0   \\
            0  &-3
         \end{mat}
         \begin{mat}[r]
            -1  &-2   \\
            -1  &-1
         \end{mat}
         =
         \begin{mat}[r]
            -5  &-4   \\
             2  &1
         \end{mat}
      \end{equation*}
      (contoh ini tidak sepenuhnya sebarang karena matriks-matriks tengah pada
      kedua ruas kiri berjumlah matriks nol).
      Perhatikan bahwa jumlah matriks-matriks serupa ini tidak serupa,
      \begin{equation*}
         \begin{mat}[r]
            1  &0   \\
            0  &3
         \end{mat}
         +
         \begin{mat}[r]
           -1  &0   \\
            0  &-3
         \end{mat}
         =
         \begin{mat}[r]
           0  &0  \\
           0  &0
         \end{mat}
         \qquad
         \begin{mat}[r]
            5/3  &-2/3   \\
            -4/3 &7/3
         \end{mat}
         +
         \begin{mat}[r]
            -5  &-4   \\
             2  &1
         \end{mat}
         \neq
         \begin{mat}[r]
           0  &0  \\
           0  &0
         \end{mat}
      \end{equation*}
      karena matriks nol hanya serupa dengan dirinya sendiri.
    \end{answer}
  \item 
    Tunjukkan bahwa jika \( T-\lambda I \) dan \( N \) merupakan matriks serupa,
    maka \( T \) dan \( N+\lambda I \) juga serupa.
    \begin{answer}
    Jika \( N=P(T-\lambda I)P^{-1} \), maka
    \( N=PTP^{-1}-P(\lambda I)P^{-1} \).
    Matriks diagonal \( \lambda I \) komutatif dengan setiap matriks, sehingga
    \( P(\lambda I)P^{-1}=PP^{-1}(\lambda I)=\lambda I \).
    Jadi, \( N=PTP^{-1}-\lambda I \), dan akibatnya
    \( N+\lambda I=PTP^{-1} \).
    (Jadi, keduanya bukan hanya serupa; keduanya bahkan serupa melalui \( P \)
    yang sama.)
    \end{answer}
\end{exercises}

















\subsection{Keterdiagonalan}
Subbagian sebelumnya menunjukkan bahwa walaupun matriks-matriks serupa pasti
ekuivalen, kebalikannya tidak berlaku.
Sebagian kelas ekuivalensi matriks terpecah menjadi dua atau lebih kelas
keserupaan; sebagai contoh, matriks nonsingular $\nbyn{2}$ membentuk satu kelas
ekuivalensi matriks tetapi lebih dari satu kelas keserupaan.

Diagram di bawah mengilustrasikannya.
Kurva utuh menunjukkan kelas ekuivalensi matriks, sedangkan pembatas putus-putus
menandai kelas keserupaan.
Setiap bintang merupakan matriks yang mewakili kelas keserupaannya.
%  so each
% similarity class has one. 
Kita tidak dapat menggunakan bentuk kanonik bagi ekuivalensi matriks, yaitu
matriks blok identitas parsial, sebagai bentuk kanonik bagi keserupaan karena
setiap kelas ekuivalensi matriks hanya memiliki satu matriks identitas parsial.
\begin{center}
  \includegraphics{jc/mp/ch5.5}
\end{center}
Untuk mengembangkan bentuk kanonik bagi wakil kelas keserupaan, wajar jika kita
membangun di atas hasil sebelumnya.
Jadi, jika suatu kelas keserupaan memang memuat matriks identitas parsial,
matriks itu harus menjadi wakil kelas tersebut.
Selain itu, setiap wakil harus sesederhana mungkin.
%(the partial identities are simple in that they consist mostly of zeros). 

Perluasan paling sederhana dari bentuk identitas parsial adalah bentuk diagonal.

\begin{definition} \label{df:Diagonalizable}
%<*df:Diagonalizable>
Suatu transformasi disebut \definend{dapat didiagonalkan}
\index{dapat didiagonalkan|(}%
\index{transformasi!dapat didiagonalkan}
jika memiliki representasi diagonal terhadap basis yang sama bagi kodomain dan
domain.
Suatu \definend{matriks yang dapat didiagonalkan}
\index{matriks!dapat didiagonalkan} adalah matriks yang serupa dengan matriks
diagonal:~\( T \) dapat didiagonalkan jika terdapat matriks nonsingular \( P \)
sedemikian sehingga \( PTP^{-1} \) diagonal.
%</df:Diagonalizable>
\end{definition}

\begin{example} \label{ex:DiagTwoByTwo}
Matriks
\begin{equation*}
  \begin{mat}[r] 
     4 &-2 \\ 
     1 &1 
  \end{mat}
\end{equation*} 
dapat didiagonalkan.
\begin{equation*}
  \begin{mat}[r]
     2  &0   \\
     0  &3
  \end{mat}
  =
  \begin{mat}[r]
    -1  &2  \\
     1  &-1
  \end{mat}
  \begin{mat}[r]
     4  &-2 \\
     1  &1
  \end{mat}
  \begin{mat}[r]
    -1  &2  \\
     1  &-1
  \end{mat}^{-1}
\end{equation*}
\end{example}

Di bawah ini kita akan melihat cara mencari matriks~$P$, tetapi terlebih dahulu
perhatikan bahwa tidak setiap matriks serupa dengan matriks diagonal.
Karena itu, bentuk diagonal tidak cukup sebagai bentuk kanonik bagi keserupaan.

\begin{example}  \label{ex:MatrixNotDiagonalizable}
%<*ex:NotDiagonalizable>
Matriks berikut tidak dapat didiagonalkan.
\begin{equation*}
  N=\begin{mat}[r]
       0  &0  \\
       1  &0
    \end{mat}
\end{equation*}
Fakta bahwa $N$ bukan matriks nol berarti bahwa $N$ tidak dapat serupa dengan
matriks nol, karena matriks nol hanya serupa dengan dirinya sendiri.
Jadi, jika $N$ serupa dengan suatu matriks diagonal~$D$, maka $D$ memiliki
sedikitnya satu entri tak nol pada diagonalnya.

Pokok pentingnya adalah bahwa suatu pangkat $N$ merupakan matriks nol; secara
khusus, $N^2$ adalah matriks nol.
Hal ini menyiratkan bahwa bagi sebarang pemetaan~$n$ yang direpresentasikan oleh
\( N \) terhadap suatu $B,B$, komposisi \( \composed{n}{n} \) merupakan
pemetaan nol.
Pada gilirannya, hal ini menyiratkan bahwa kuadrat sebarang matriks yang
merepresentasikan~$n$ terhadap suatu $\hat{B},\hat{B}$ merupakan matriks nol.
Namun, bagi sebarang matriks diagonal tak nol~$D$, entri-entri $D^2$ adalah
kuadrat entri-entri~$D$, sehingga $D^2$ tidak mungkin merupakan matriks nol.
Jadi, $N$ tidak dapat didiagonalkan.
%</ex:NotDiagonalizable>
\end{example}

Jadi, tidak setiap kelas keserupaan memuat matriks diagonal.
Sekarang kita mencirikan kapan suatu matriks dapat didiagonalkan.
% However, some similarity classes do contain a diagonal matrix and 
% the canonical form that we are developing has the property that if
% a matrix can be diagonalized then the diagonal matrix is the canonical
% representative of its similarity class. 

\begin{lemma} \label{lm:DiagIffBasisOfEigens}
%<*lm:DiagIffBasisOfEigens>
Suatu transformasi \( t \) dapat didiagonalkan jika dan hanya jika terdapat
basis
\( B=\sequence{\vec{\beta}_1,\ldots,\vec{\beta}_n } \)
dan skalar \( \lambda_1,\ldots,\lambda_n \) sedemikian sehingga
\( t(\vec{\beta}_i)=\lambda_i\vec{\beta}_i \)
bagi setiap \( i \).
%</lm:DiagIffBasisOfEigens>
\end{lemma}

\begin{proof}
%<*pf:DiagIffBasisOfEigens>
Tinjau suatu matriks representasi diagonal.
\begin{equation*}
   \rep{t}{B,B}=
   \begin{pmat}{c@{\hspace*{1em}}c@{\hspace*{1em}}c}
      \vdots                    &       &\vdots                     \\
      \rep{t(\vec{\beta}_1)}{B} &\cdots &\rep{t(\vec{\beta}_n)}{B}  \\
      \vdots                    &       &\vdots
   \end{pmat}
   =
   \begin{pmat}{c@{\hspace*{1em}}c@{\hspace*{1em}}c}
      \lambda_1   &       &0         \\
      \vdots      &\ddots &\vdots    \\
      0           &       &\lambda_n
   \end{pmat}
\end{equation*}
Tinjau representasi suatu anggota basis ini terhadap basis tersebut,
$\rep{\vec{\beta}_i}{B}$.
Hasil kali matriks diagonal dan vektor representasi
\begin{equation*}
   \rep{t(\vec{\beta}_i)}{B}
   =\begin{pmat}{c@{\hspace*{1em}}c@{\hspace*{1em}}c}
      \lambda_1   &       &0         \\
      \vdots      &\ddots &\vdots    \\
      0           &       &\lambda_n
   \end{pmat}  
  \colvec{0 \\ \vdots \\ 1 \\ \vdots \\ 0}
  =\colvec{0 \\ \vdots \\ \lambda_i \\ \vdots \\ 0}
\end{equation*}
memiliki tindakan yang dinyatakan.
%</pf:DiagIffBasisOfEigens>
\end{proof}

\begin{example}     \label{ex:DiagUpperTrian}
Untuk mendiagonalkan
\begin{equation*}
   T=\begin{mat}[r]
      3  &2  \\
      0  &1
   \end{mat}
\end{equation*}
ambil $T$ sebagai representasi suatu transformasi terhadap basis standar,
$\rep{t}{\stdbasis_2,\stdbasis_2}$, lalu cari basis
\( B=\sequence{\vec{\beta}_1,\vec{\beta}_2} \) sedemikian sehingga
\begin{equation*}
  \rep{t}{B,B}
  =
  \begin{mat}
    \lambda_1  &0          \\
    0          &\lambda_2
  \end{mat}
\end{equation*}
yakni sedemikian sehingga
$t(\vec{\beta}_1)=\lambda_1\vec{\beta}_1$ dan
$t(\vec{\beta}_2)=\lambda_2\vec{\beta}_2$.
\begin{equation*}
  \begin{mat}[r]
     3  &2  \\
     0  &1
  \end{mat}
  \vec{\beta}_1=\lambda_1\cdot\vec{\beta}_1
  \qquad
  \begin{mat}[r]
     3  &2  \\
     0  &1
  \end{mat}
  \vec{\beta}_2=\lambda_2\cdot\vec{\beta}_2
\end{equation*} 
Kita mencari skalar \( x \) sedemikian sehingga persamaan
\begin{equation*}
 \begin{mat}[r]
    3  &2  \\
    0  &1
 \end{mat}
 \colvec{b_1 \\ b_2}=x\cdot\colvec{b_1 \\ b_2}
\end{equation*}
memiliki solusi $b_1$ dan $b_2$ yang tidak keduanya bernilai $0$ (vektor nol
bukan anggota sebarang basis).
Ini merupakan sistem linear.
\begin{equation*}
  \begin{linsys}{2}
     (3-x)\cdot b_1  &+  &2\cdot b_2       &=  &0  \\
                     &   &(1-x)\cdot b_2   &=  &0 
  \end{linsys}
\tag*{($*$)}\end{equation*}
Perhatikan dahulu persamaan terbawah.
Ada dua kasus:~$b_2=0$ atau~$x=1$.

Dalam kasus \( b_2=0 \), persamaan pertama memberikan $b_1=0$ atau \( x=3 \).
Karena kemungkinan $b_2=0$ dan~$b_1=0$ sekaligus telah dilarang, tersisa entri
diagonal pertama $\lambda_1=3$.
Dengan nilai tersebut, persamaan pertama pada ($*$) menjadi
$0\cdot b_1+2\cdot b_2=0$. Jadi, vektor-vektor yang terkait dengan
$\lambda_1=3$ memiliki komponen kedua nol, sedangkan komponen pertama bebas.
\begin{equation*}
     \begin{mat}[r]
        3  &2  \\
        0  &1
     \end{mat}
     \colvec{b_1 \\ 0}=3\cdot\colvec{b_1 \\ 0} 
\end{equation*}
Untuk memperoleh vektor basis pertama, pilih sebarang $b_1$ tak nol.
\begin{equation*}
   \vec{\beta}_1=\colvec[r]{1 \\ 0}
\end{equation*}

Kasus lain bagi persamaan terbawah pada ($*$) adalah $\lambda_2=1$.
Persamaan pertama pada ($*$) kemudian menjadi
$2\cdot b_1+2\cdot b_2=0$. Jadi, vektor-vektor yang terkait dengan kasus ini
memiliki komponen kedua yang merupakan negatif komponen pertama.
\begin{equation*}
     \begin{mat}[r]
        3  &2  \\
        0  &1
     \end{mat}
     \colvec{b_1 \\ -b_1}=1\cdot\colvec{b_1 \\ -b_1} 
\end{equation*}
Peroleh vektor basis kedua dengan memilih salah satu vektor tak nol tersebut.
\begin{equation*}
   \vec{\beta}_2=\colvec[r]{1 \\ -1}
\end{equation*}

Sekarang gambarkan diagram keserupaan (ingat bahwa kita bekerja dengan skalar
kompleks, bukan real),
\begin{equation*}
     \begin{CD}
            \C^2_{\wrt{\stdbasis_2}}                   
               @>t>T>        
               \C^2_{\wrt{\stdbasis_2}}       \\
            @V{\scriptstyle\identity} VV                
               @V{\scriptstyle\identity} VV \\
            \C^2_{\wrt{B}}         
               @>t>D>        
               \C^2_{\wrt{B}}
      \end{CD}
\end{equation*}
dan perhatikan bahwa matriks $\rep{\identity}{B,\stdbasis_2}$ mudah diperoleh,
sehingga memberikan diagonalisasi berikut.
\begin{equation*}
   \begin{mat}[r]
     3  &0  \\
     0  &1
   \end{mat}
   =
   \begin{mat}[r]
     1  &1  \\
     0  &-1
   \end{mat}^{-1}
   \begin{mat}[r]
     3  &2  \\
     0  &1
   \end{mat}
   \begin{mat}[r]
     1  &1  \\
     0  &-1
   \end{mat}
\end{equation*}
\end{example}

Sisa bagian ini mengembangkan contoh tersebut dengan meninjau lebih dekat sifat
pada \nearbylemma{lm:DiagIffBasisOfEigens}, termasuk cara yang lebih ringkas
untuk mencari nilai-nilai $\lambda$.
Bagian setelahnya mengembangkan \nearbyexample{ex:MatrixNotDiagonalizable} untuk
memahami apa yang dapat menghalangi diagonalisasi.
Kemudian, bagian terakhir menggabungkan keduanya untuk menghasilkan bentuk
kanonik yang, dalam suatu pengertian, sesederhana mungkin.

\begin{exercises}
  \recommended \item 
    Ulangi \nearbyexample{ex:DiagUpperTrian} bagi matriks dari
    \nearbyexample{ex:DiagTwoByTwo}.
    \begin{answer}
      Karena kita memilih vektor-vektor basis secara sebarang, banyak jawaban
      berbeda dapat diperoleh.
      Namun, berikut salah satu caranya.
      Untuk mendiagonalkan
      \begin{equation*}
         T=\begin{mat}[r]
            4  &-2  \\
            1  &1
         \end{mat}
      \end{equation*}
      ambil matriks itu sebagai representasi suatu transformasi terhadap basis
      standar, $T=\rep{t}{\stdbasis_2,\stdbasis_2}$, lalu cari
      \( B=\sequence{\vec{\beta}_1,\vec{\beta}_2} \) sedemikian sehingga
      \begin{equation*}
        \rep{t}{B,B}
        =
        \begin{mat}
          \lambda_1  &0          \\
          0          &\lambda_2
        \end{mat}
      \end{equation*}
      yakni sedemikian sehingga
      $t(\vec{\beta}_1)=\lambda_1\vec{\beta}_1$ dan
      $t(\vec{\beta}_2)=\lambda_2\vec{\beta}_2$.
      \begin{equation*}
        \begin{mat}[r]
           4  &-2  \\
           1  &1
        \end{mat}
        \vec{\beta}_1=\lambda_1\cdot\vec{\beta}_1
        \qquad
        \begin{mat}[r]
           4  &-2  \\
           1  &1
        \end{mat}
        \vec{\beta}_2=\lambda_2\cdot\vec{\beta}_2
      \end{equation*} 
      Kita mencari skalar \( x \) sedemikian sehingga persamaan
      \begin{equation*}
       \begin{mat}[r]
          4  &-2  \\
          1  &1
       \end{mat}
       \colvec{b_1 \\ b_2}=x\cdot\colvec{b_1 \\ b_2}
      \end{equation*}
      memiliki solusi $b_1$ dan $b_2$ yang tidak keduanya nol.
      Tulis ulang sebagai sistem linear,
      \begin{equation*}
        \begin{linsys}{2}
           (4-x)\cdot b_1  &+  &-2\cdot b_2       &=  &0  \\
           1\cdot b_1      &+   &(1-x)\cdot b_2   &=  &0 
        \end{linsys}
      \end{equation*}
      Jika $x=4$, persamaan pertama memberikan $b_2=0$, lalu persamaan kedua
      memberikan $b_1=0$.
      Kasus ketika kedua $b$ bernilai nol telah dilarang, sehingga kita dapat
      mengandaikan $x\neq 4$.
      \begin{equation*}
        \grstep{(-1/(4-x))\rho_1+\rho_2}
        \begin{linsys}{2}
           (4-x)\cdot b_1  &+   &-2\cdot b_2                   &=  &0  \\
                           &    &((x^2-5x+6)/(4-x))\cdot b_2   &=  &0 
        \end{linsys} 
      \end{equation*}
      Tinjau persamaan terbawah.
      Jika \( b_2=0 \), persamaan pertama memberikan $b_1=0$ atau $x=4$.
      Kasus $b_1=b_2=0$ tidak diizinkan.
      Kemungkinan lain bagi persamaan terbawah adalah pembilang pecahan
      $x^2-5x+6=(x-2)(x-3)$ bernilai nol.
      Kasus $x=2$ memberikan persamaan pertama $2b_1-2b_2=0$. Jadi, dengan
      $x=2$ terkait vektor-vektor yang komponen pertama dan keduanya sama:
      \begin{equation*}
         \vec{\beta}_1=\colvec[r]{1 \\ 1}
         \qquad\text{(sehingga \(
           \begin{mat}[r]
              4  &-2  \\
              1  &1
           \end{mat}
           \colvec[r]{1 \\ 1}=2\cdot\colvec[r]{1 \\ 1} \), dan $\lambda_1=2$).}
      \end{equation*}
      Jika \( x=3 \), persamaan pertama adalah $b_1-2b_2=0$. Jadi, vektor yang
      terkait memiliki komponen pertama dua kali komponen keduanya:
      \begin{equation*}
         \vec{\beta}_2=\colvec[r]{2 \\ 1}
         \qquad\text{(sehingga \(
           \begin{mat}[r]
              4  &-2  \\
              1  &1
           \end{mat}
           \colvec[r]{2 \\ 1}=3\cdot\colvec[r]{2 \\ 1} \), sehingga $\lambda_2=3$).}
      \end{equation*}
      Diagram berikut
        \begin{equation*}
          \begin{CD}
            \C^2_{\wrt{\stdbasis_2}}                   
               @>t>T>        
               \C^2_{\wrt{\stdbasis_2}}       \\
            @V{\scriptstyle\identity} VV                
               @V{\scriptstyle\identity} VV \\
            \C^2_{\wrt{B}}         
               @>t>D>        
               \C^2_{\wrt{B}}
          \end{CD}
        \end{equation*}
      menunjukkan cara memperoleh diagonalisasinya.
      \begin{equation*}
         \begin{mat}[r]
           2  &0  \\
           0  &3
         \end{mat}
         =
         \begin{mat}[r]
           1  &2  \\
           1  &1
         \end{mat}^{-1}
         \begin{mat}[r]
           4  &-2  \\
           1  &1
         \end{mat}
         \begin{mat}[r]
           1  &2  \\
           1  &1
         \end{mat}
      \end{equation*}
      \textit{Komentar}.
      Dengan penggantian nama berikut, persamaan ini sesuai dengan definisi
      $T=PSP^{-1}$.
      \begin{equation*}
        T=
         \begin{mat}[r]
           2  &0  \\
           0  &3
         \end{mat}
         \quad
         P=
         \begin{mat}[r]
           1  &2  \\
           1  &1
         \end{mat}^{-1}
         \quad
         P^{-1}=
         \begin{mat}[r]
           1  &2  \\
           1  &1
         \end{mat}
         \quad
         S=
         \begin{mat}[r]
           4  &-2  \\
           1  &1
         \end{mat}
      \end{equation*}
    \end{answer}
  \recommended \item \label{exer:DiagTheseA} 
    Diagonalkan matriks berikut dengan mengikuti langkah-langkah pada
    \nearbyexample{ex:DiagUpperTrian}.
    \begin{equation*}
      \begin{mat}[r]
      1  &1  \\
      0  &0
      \end{mat}
    \end{equation*}
    \begin{exparts}
      \partsitem Susun persamaan matriks-vektor yang dideskripsikan dalam
         \nearbylemma{lm:DiagIffBasisOfEigens}, lalu tulis ulang sebagai sistem
         linear.
      \partsitem Dengan meninjau solusi sistem tersebut, cari dua vektor untuk
         membentuk basis.
         (Tinjau secara terpisah kasus $x=0$ dan $x\neq 0$.
         Ingat pula bahwa vektor nol tidak dapat menjadi anggota basis.)
      \partsitem Gunakan basis tersebut dalam diagram keserupaan untuk
        memperoleh matriks diagonal sebagai hasil kali tiga matriks lain.
    \end{exparts}
    \begin{answer}
     \begin{exparts}
        \partsitem Kita menginginkan basis yang vektor-vektor anggotanya memenuhi
          \begin{equation*}
            \begin{mat}
              1 &1 \\
              0 &0
            \end{mat}
            \colvec{b_1 \\ b_2}
            =x\cdot \colvec{b_1 \\ b_2}
          \end{equation*}
          sehingga kita meninjau sistem linear yang dihasilkan.
          \begin{equation*}
            \begin{linsys}{2}
              b_1 &+ &b_2 &= &x\cdot b_1 \\
                  &  &0   &= &x\cdot b_2 \\
            \end{linsys}
          \end{equation*}
        \partsitem
          Jika $x=0$, maka $b_1+b_2=0$, sehingga $b_1=-b_2$.
          Artinya, dengan kasus ini terkait unsur-unsur basis yang entrinya
          memiliki tanda berlawanan.
          (Entri tersebut harus tak nol karena vektor nol tidak dapat menjadi
          anggota basis.)

          Jika $x\neq 0$, maka $b_2=0$, dan substitusi ke persamaan teratas
          memberikan $b_1=x\cdot b_1$.
          Jadi, dalam kasus ini $b_1=0$ atau $x=1$.
          Namun, karena $b_2=0$ dalam kasus ini, $b_1=0$ tidak mungkin sebab
          vektor nol tidak termuat dalam sebarang basis.
          Jadi, $x=1$ dan sistemnya menjadi
          \begin{equation*}
            \begin{linsys}{2}
              b_1 &+ &b_2 &= &1\cdot b_1 \\
                  &  &0   &= &1\cdot b_2 \\
            \end{linsys}
          \end{equation*}
          Jadi, dengan kasus ini terkait unsur-unsur basis yang memenuhi
          $b_2=0$, sedangkan~$b_1$ dapat berupa sebarang nilai tak nol.
           
          Karena itu, salah satu basis yang sesuai adalah
          \begin{equation*}
            B=\sequence{\colvec{1 \\ -1},
                        \colvec{1 \\ 0}}
          \end{equation*}
        \partsitem
          Dengan basis tersebut, gambarkan diagram keserupaan untuk memperoleh
          diagonalisasi.
          \begin{equation*}
            \begin{CD}
            \C^2_{\wrt{\stdbasis_2}}                   
               @>t>T>        
               \C^2_{\wrt{\stdbasis_2}}       \\
            @V{\scriptstyle\identity} VV                
               @V{\scriptstyle\identity} VV \\
            \C^2_{\wrt{B}}         
               @>t>D>        
               \C^2_{\wrt{B}}
           \end{CD}
          \end{equation*}
          Dalam diagram ini,
          \begin{equation*}
             T=
             \begin{mat}
               1  &1  \\
               0  &0
             \end{mat}
          \end{equation*}
          dan kita akan menghitung matriks~$D$ menggunakan basis~$B$ dengan
          menelusuri jalur dari kiri bawah ke kiri atas, lalu ke kanan atas, dan
          akhirnya ke kanan bawah.

          Mula-mula kita bergerak dari kiri bawah ke kiri atas.
          Karena basis pada bagian atas diagram dipilih sebagai basis standar
          $\stdbasis_2$, matriks yang merepresentasikan pemetaan itu mudah
          diperoleh.
          \begin{equation*}
            \rep{\identity}{B,\stdbasis_2}=
            \begin{mat}
              1 &1 \\
             -1 &0
            \end{mat}
          \end{equation*}
          Perhitungan bagi seluruh jalur memang memberikan hasil diagonal
          (tentu saja, bagian kiri bawah ke kiri atas muncul pada sisi kanan
          hasil kali ketiga matriks).
          \begin{align*}
                 D
                 &=
                 \begin{mat}
                    1  &1  \\
                    -1  &0
                 \end{mat}^{-1}
                 \begin{mat}
                    1  &1  \\
                    0  &0
                 \end{mat}
                 \begin{mat}
                    1  &1  \\
                    -1  &0
                 \end{mat}                \\
                 &=
                 \begin{mat}
                    0  &-1  \\
                    1  &1
                 \end{mat}
                 \begin{mat}
                    1  &1  \\
                    0  &0
                 \end{mat}
                 \begin{mat}
                    1  &1  \\
                    -1  &0
                 \end{mat}                
                =\begin{mat}
                    0  &0  \\
                    0  &1
                 \end{mat}
          \end{align*}
     \end{exparts}    
    \end{answer}
  \recommended \item \label{exer:DiagTheseB} 
    Ikuti \nearbyexample{ex:DiagUpperTrian} untuk mendiagonalkan matriks ini.
    \begin{equation*}
      \begin{mat}[r]
      0  &1  \\
      1  &0
      \end{mat}
    \end{equation*}
    \begin{exparts}
      \partsitem Susun persamaan matriks-vektor, lalu tulis ulang sebagai sistem
         linear.
      \partsitem Dengan meninjau solusi sistem tersebut, cari dua vektor untuk
         membentuk basis.
         (Tinjau secara terpisah kasus $x=0$ dan $x\neq 0$.
         Ingat pula bahwa vektor nol tidak dapat menjadi anggota basis.)
      \partsitem Dengan basis tersebut, gunakan diagram keserupaan untuk
         memperoleh diagonalisasi sebagai hasil kali tiga matriks.
    \end{exparts}
    \begin{answer}
     \begin{exparts}
        \partsitem 
          Anggota-anggota basis memenuhi
          \begin{equation*}
            \begin{mat}
              0 &1 \\
              1 &0
            \end{mat}
            \colvec{b_1 \\ b_2}
            =x\cdot \colvec{b_1 \\ b_2}
          \end{equation*}
          sehingga sistem linear yang dihasilkan adalah
          \begin{equation*}
            \begin{linsys}{2}
                   &  &b_2 &= &x\cdot b_1 \\
               b_1 &  &    &= &x\cdot b_2 \\
            \end{linsys}
            \qquad\Longrightarrow\qquad
            \begin{linsys}{2}
               -xb_1 &+ &b_2  &= &0 \\
                 b_1 &- &xb_2 &= &0 \\
            \end{linsys}
          \end{equation*}
        \partsitem
          Jika $x=0$, maka $b_1=b_2=0$. Namun, vektor nol bukan anggota sebarang
          basis, sehingga kasus ini tidak mungkin terjadi.

          Kasus lainnya adalah $x\neq 0$.
          \begin{equation*}
            \begin{linsys}{2}
               -xb_1 &+ &b_2  &= &0 \\
                 b_1 &- &xb_2 &= &0 \\
            \end{linsys}  
            \grstep{(1/x)\rho_1+\rho_2}
            \begin{linsys}{2}
               -xb_1 &+ &b_2           &= &0 \\
                     &  &(-x+(1/x))b_2 &= &0 \\
            \end{linsys}  
          \end{equation*}
          Tinjau persamaan terbawah.
          Kemungkinannya adalah $b_2=0$ atau $-x+(1/x)=0$.
          Jika $b_2=0$, substitusi ke persamaan teratas memberikan $b_1=0$
          karena asumsi kasus ini adalah $x\neq 0$.
          Namun, vektor nol bukan anggota sebarang basis, sehingga hal ini tidak
          mungkin terjadi.

          Tersisa
          \begin{equation*}
            \frac{1}{x}-x=0
            \qquad\Longrightarrow\qquad
            0=\frac{1-x^2}{x}=\frac{(1-x)(1+x)}{x}
          \end{equation*}
          Jika $x=1$, sistemnya adalah
          \begin{equation*}
            \begin{linsys}{2}
               -1\cdot b_1 &+ &b_2  &= &0 \\
                           &  &0    &= &0 \\
            \end{linsys}
          \end{equation*}
          sehingga vektor-vektor basis memenuhi $b_2=b_1$.
          Jika $x=-1$, sistemnya adalah
          \begin{equation*}
            \begin{linsys}{2}
               b_1 &+ &b_2  &= &0 \\
                   &  &0    &= &0 \\
            \end{linsys}
          \end{equation*}
          dan vektor-vektor basis memenuhi $b_2=-b_1$.
          Karena itu, berikut salah satu contoh basis yang sesuai.
          \begin{equation*}
            B=\sequence{\colvec{1 \\ 1},
                        \colvec{1 \\ -1}}
          \end{equation*}
        \partsitem
          Untuk memperoleh diagonalisasi, gambarkan diagram keserupaan
          \begin{equation*}
            \begin{CD}
            \C^2_{\wrt{\stdbasis_2}}                   
               @>t>T>        
               \C^2_{\wrt{\stdbasis_2}}       \\
            @V{\scriptstyle\identity} VV                
               @V{\scriptstyle\identity} VV \\
            \C^2_{\wrt{B}}         
               @>t>D>        
               \C^2_{\wrt{B}}
           \end{CD}
          \end{equation*}
          dengan $T$ sebagai matriks yang diberikan dalam soal.
          \begin{equation*}
             T=
             \begin{mat}
               0  &1  \\
               1  &0
             \end{mat}
          \end{equation*}
          Hitung matriks~$D$ menggunakan basis~$B$ dengan bergerak dari kiri
          bawah ke kiri atas, lalu ke kanan atas, dan akhirnya ke kanan bawah.

          Perhatikan bahwa matriks yang merepresentasikan pemetaan dari kiri
          bawah ke kiri atas mudah diperoleh.
          \begin{equation*}
            \rep{\identity}{B,\stdbasis_2}=
            \begin{mat}
              1 &1 \\
              1 &-1
            \end{mat}
          \end{equation*}
          Dengan itu, perhitungan bagi seluruh jalur memang memberikan hasil
          diagonal.
          \begin{align*}
                 D
                 &=
                 \begin{mat}
                    1  &1  \\
                    1  &-1
                 \end{mat}^{-1}
                 \begin{mat}
                    0  &1  \\
                    1  &0
                 \end{mat}
                 \begin{mat}
                    1  &1  \\
                    1  &-1
                 \end{mat}                \\
                 &=
                 -\frac{1}{2}
                 \begin{mat}
                    -1  &-1  \\
                    -1  &1
                 \end{mat}
                 \begin{mat}
                    0  &1  \\
                    1  &0
                 \end{mat}
                 \begin{mat}
                    1  &1  \\
                    1  &-1
                 \end{mat}                
                =\begin{mat}
                    1  &0  \\
                    0  &-1
                 \end{mat}
          \end{align*}
      \end{exparts}  
    \end{answer}
  \item 
    Diagonalkan matriks-matriks segitiga atas berikut.
    \begin{exparts*}
      \partsitem 
        $\begin{mat}[r]
          -2  &1  \\
           0  &2
        \end{mat}$
      \partsitem 
        $\begin{mat}[r]
          5  &4  \\
          0  &1
        \end{mat}$
    \end{exparts*}
    \begin{answer}
      \begin{exparts}
        \partsitem
          Menyusun
          \begin{equation*}
            \begin{mat}[r]
              -2  &1  \\
               0  &2
            \end{mat}
            \colvec{b_1 \\ b_2}
            =x\cdot\colvec{b_1 \\ b_2}
            \qquad\Longrightarrow\qquad 
            \begin{linsys}{2}
               (-2-x)\cdot b_1  &+  &b_2            &=  &0  \\
                                &   &(2-x)\cdot b_2 &= &0
            \end{linsys}
          \end{equation*}
          memberikan dua kemungkinan, yaitu $b_2=0$ atau $x=2$.
          Mengikuti kemungkinan $b_2=0$ menghasilkan persamaan pertama
          $(-2-x)b_1=0$, dengan dua kasus $b_1=0$ atau $x=-2$.
          Jadi, pada kemungkinan pertama ini kita memperoleh $x=-2$ dan vektor-
          vektor terkait yang komponen keduanya nol, sedangkan komponen
          pertamanya bebas.
          \begin{equation*}
            \begin{mat}[r]
              -2  &1  \\
               0  &2
            \end{mat}
            \colvec{b_1 \\ 0}
            =-2\cdot\colvec{b_1 \\ 0}
            \qquad
            \vec{\beta}_1=\colvec[r]{1 \\ 0}
          \end{equation*}
          Mengikuti kemungkinan lainnya menghasilkan persamaan pertama
          $-4b_1+b_2=0$. Jadi, vektor-vektor yang terkait dengan solusi ini
          memiliki komponen kedua empat kali komponen pertamanya.
          \begin{equation*}
            \begin{mat}[r]
              -2  &1  \\
               0  &2
            \end{mat}
            \colvec{b_1 \\ 4b_1}
            =2\cdot\colvec{b_1 \\ 4b_1}
            \qquad
            \vec{\beta}_2=\colvec{1 \\ 4}
          \end{equation*}
          Diagonalisasinya adalah
          \begin{equation*}
            \begin{mat}[r]
              1  &1  \\
              0  &4
            \end{mat}^{-1}
            \begin{mat}[r]
              -2  &1  \\
               0  &2
            \end{mat}
            \begin{mat}[r]
              1  &1  \\
              0  &4
            \end{mat}
            =
            \begin{mat}[r]
              -2  &0  \\
               0  &2
            \end{mat}
          \end{equation*}
      \partsitem Perhitungannya serupa dengan bagian sebelumnya.
          \begin{equation*}
            \begin{mat}[r]
               5  &4  \\
               0  &1
            \end{mat}
            \colvec{b_1 \\ b_2}
            =x\cdot\colvec{b_1 \\ b_2}
            \qquad\Longrightarrow\qquad 
            \begin{linsys}{2}
               (5-x)\cdot b_1  &+  &4\cdot b_2     &=  &0  \\
                               &   &(1-x)\cdot b_2 &=  &0
            \end{linsys}
          \end{equation*}
          Persamaan terbawah memberikan dua kemungkinan, yaitu $b_2=0$ atau
          $x=1$.
          Mengikuti kemungkinan $b_2=0$ dan menyingkirkan kasus ketika $b_2$
          dan $b_1$ sama-sama nol memberikan $x=5$, yang terkait dengan vektor-
          vektor yang komponen keduanya nol dan komponen pertamanya bebas.
          \begin{equation*}
            \vec{\beta}_1=\colvec[r]{1 \\ 0}
          \end{equation*}
          Kemungkinan $x=1$ memberikan persamaan pertama $4b_1+4b_2=0$.
          Jadi, vektor-vektor terkait memiliki komponen kedua yang merupakan
          negatif komponen pertamanya.
          \begin{equation*}
            \vec{\beta}_2=\colvec[r]{1 \\ -1}
          \end{equation*}
          Dengan demikian, diperoleh diagonalisasi berikut.
          \begin{equation*}
            \begin{mat}[r]
              1  &1  \\
              0  &-1
            \end{mat}
            \begin{mat}[r]
               5  &4  \\
               0  &1
            \end{mat}
            \begin{mat}[r]
              1  &1  \\
              0  &-1
            \end{mat}^{-1}
            =
            \begin{mat}[r]
               5  &0  \\
               0  &1
            \end{mat}
          \end{equation*}
      \end{exparts}
    \end{answer}
  \item Jika kita mencoba mendiagonalkan matriks pada
    \nearbyexample{ex:MatrixNotDiagonalizable},
    \begin{equation*}
       N=\begin{mat}[r]
            0  &0  \\
            1  &0
         \end{mat}
    \end{equation*}
    menggunakan metode pada \nearbyexample{ex:DiagUpperTrian}, apa yang gagal?
    \begin{exparts}
      \partsitem Gambarkan diagram keserupaan dengan $N$.
      \partsitem Susun persamaan matriks-vektor yang dideskripsikan dalam
         \nearbylemma{lm:DiagIffBasisOfEigens}, lalu tulis ulang sebagai sistem
         linear.
      \partsitem Dengan meninjau solusi sistem tersebut, temukan masalahnya.
         (Tinjau secara terpisah kasus $x=0$ dan $x\neq 0$.)
    \end{exparts}
    \begin{answer}
      \begin{exparts}
        \item Dalam diagram keserupaan berikut,
          \begin{equation*}
            \begin{CD}
            \C^2_{\wrt{\stdbasis_2}}                   
               @>n>N>        
               \C^2_{\wrt{\stdbasis_2}}       \\
            @V{\scriptstyle\identity} VV                
               @V{\scriptstyle\identity} VV \\
            \C^2_{\wrt{B}}         
               @>n>D>        
               \C^2_{\wrt{B}}
           \end{CD}
          \end{equation*}
          kita akan mencoba menghitung matriks~$D$ dan mencari basis~$B$ yang
          sesuai dengan menelusuri jalur dari kiri bawah ke kiri atas, lalu ke
          kanan atas, dan akhirnya ke kanan bawah.
        \item Diagram menunjukkan bahwa bagi
          $B=\sequence{\vec{\beta}_1,\vec{\beta}_2}$, kita menginginkan
          $n(\vec{\beta}_1)=\lambda_1\vec{\beta}_1$ dan
          $n(\vec{\beta}_2)=\lambda_2\vec{\beta}_2$.
          Untuk menentukan $\lambda_1$ dan $\lambda_2$, ketika matriks~$N$
          merepresentasikan pemetaan terhadap basis standar, kita ingin mencari
          solusi sistem berikut.
          \begin{equation*}
            \begin{mat}
              0 &0 \\
              1 &0
            \end{mat}
            \colvec{b_1 \\ b_2}
            =
            x\cdot\colvec{b_1 \\ b_2}           
          \end{equation*}
          Sistem linear ekuivalennya
          \begin{equation*}
            \begin{linsys}{2}
              0b_1 &+  &0b_2 &= &xb_1 \\
              b_1  &+  &0b_2 &= &xb_2 \\
            \end{linsys}
            \qquad
            \Longrightarrow
            \qquad
            \begin{linsys}{2}
              -xb_1 &   &     &= &0 \\
              b_1   &-  &xb_2 &= &0 \\
            \end{linsys}
          \end{equation*}
          Menurut baris pertama, $x=0$ atau~$b_1=0$.
          Jika $x=0$, persamaan kedua menjadi $b_1-0=0$, sehingga $b_1=0$ dan
          kita memperoleh himpunan solusi berikut.
          \begin{equation*}
            \set{\colvec{0 \\ 1}b_2 \suchthat b_2\in\C}
          \end{equation*}
          Kita dapat mengambil anggota pertama basis sebagai
          $\vec{\beta}_1=\vec{e}_2$.

          Jika $x\neq 0$, persamaan pertama sistem linear memberikan $b_1=0$.
          Baris pertama sistem menjadi $0=0$, sedangkan baris kedua menjadi
          $-xb_2=0$.
          Karena kasus ini mengambil $x\neq 0$, kita memperoleh $b_2=0$ dan
          himpunan solusinya trivial.
          \begin{equation*}
            \set{\colvec{0 \\ 0}}
          \end{equation*}
          Jadi, masalahnya adalah bahwa kita tidak memperoleh basis yang
          disyaratkan oleh \nearbylemma{lm:DiagIffBasisOfEigens}.
      \end{exparts}
    \end{answer}
  \recommended \item \label{exer:PowersOfDiags}
    Apa bentuk pangkat-pangkat suatu matriks diagonal?
    \begin{answer}
      Bagi setiap bilangan bulat tak negatif \( p \), berlaku
      \begin{equation*}
        \begin{mat}
          d_1  &0      &   \\
          0    &\ddots &   \\
               &       &d_n
        \end{mat}^p=
        \begin{mat}
          d_1^p  &0      &   \\
          0      &\ddots &   \\
                 &       &d_n^p
        \end{mat}
     \end{equation*}
     Rumus yang sama berlaku bagi \( p<0 \) jika dan hanya jika semua
     $d_i\neq0$, sehingga matriks diagonal tersebut dapat diinverskan.
    \end{answer}
  \item 
     Berikan dua matriks diagonal berukuran sama yang tidak serupa.
     Apakah sebarang dua matriks diagonal berbeda harus berasal dari kelas
     keserupaan yang berbeda?
     \begin{answer}
       Kedua matriks berikut tidak serupa,
       \begin{equation*}
          \begin{mat}[r]
             0  &0  \\
             0  &0
          \end{mat}
          \qquad
          \begin{mat}[r]
             1  &0  \\
             0  &1
          \end{mat}
       \end{equation*}
       karena masing-masing merupakan satu-satunya anggota kelas keserupaannya.

       Untuk pertanyaan kedua, matriks-matriks berikut
       \begin{equation*}
          \begin{mat}[r]
             2  &0  \\
             0  &3
          \end{mat}
          \qquad
          \begin{mat}[r]
             3  &0  \\
             0  &2
          \end{mat}
       \end{equation*}
       serupa melalui matriks yang mengubah basis dari
       \( \sequence{\vec{\beta}_1,\vec{\beta}_2} \) menjadi
       \( \sequence{\vec{\beta}_2,\vec{\beta}_1} \).
       (\textit{Pertanyaan.}
        Apakah dua matriks diagonal serupa jika dan hanya jika entri-entri
        diagonalnya merupakan permutasi satu sama lain?)
    \end{answer}
  \item 
    Berikan suatu matriks diagonal nonsingular.
    Dapatkah matriks diagonal bersifat singular?
    \begin{answer}
      Bandingkan kedua matriks berikut.
      \begin{equation*}
         \begin{mat}[r]
           2  &0  \\
           0  &1
         \end{mat}
         \qquad
         \begin{mat}[r]
           2  &0  \\
           0  &0
         \end{mat}
      \end{equation*}
      Matriks pertama nonsingular, sedangkan matriks kedua singular.
     \end{answer}
  \recommended \item
    Tunjukkan bahwa invers suatu matriks diagonal adalah matriks diagonal dari
    invers setiap entrinya, jika tidak ada entri diagonal yang nol.
    Apa yang terjadi ketika suatu entri diagonal bernilai nol?
    \begin{answer}  
       Untuk memeriksa bahwa invers matriks diagonal adalah matriks diagonal
       dari invers setiap entrinya, kalikan saja.
       \begin{equation*}
          \begin{mat}
             a_{1,1}  &0                \\
             0        &a_{2,2}          \\
                      &       &\ddots    \\
                      &       &      &a_{n,n}
          \end{mat}
           \begin{mat}
             1/a_{1,1}  &0                \\
              0        &1/a_{2,2}          \\
                       &       &\ddots    \\
                       &       &      &1/a_{n,n}
           \end{mat}
           =I
      \end{equation*}
      (Menunjukkan bahwa matriks itu merupakan invers kiri sama mudahnya.)

      Jika suatu entri diagonal nol, matriks diagonal tersebut singular;
      determinannya nol.
    \end{answer}
  \item 
    Persamaan pada akhir \nearbyexample{ex:DiagUpperTrian},
    \begin{equation*}
       \begin{mat}[r]
         1  &1  \\
         0  &-1
       \end{mat}^{-1}
       \begin{mat}[r]
         3  &2  \\
         0  &1
       \end{mat}
       \begin{mat}[r]
         1  &1  \\
         0  &-1
       \end{mat}
       =
       \begin{mat}[r]
         3  &0  \\
         0  &1
       \end{mat}
    \end{equation*}
    tampak agak mengusik karena sebagai $P$ kita harus mengambil matriks
    pertama, yang ditampilkan sebagai suatu invers, sedangkan sebagai $P^{-1}$
    kita mengambil invers matriks pertama. Akibatnya, kedua pangkat $-1$ saling
    meniadakan dan matriks ini ditampilkan tanpa superskrip $-1$.
    \begin{exparts}
      \partsitem Periksa bahwa persamaan yang tampak lebih rapi berikut berlaku.
        \begin{equation*}
           \begin{mat}[r]
             3  &0  \\
             0  &1
           \end{mat}
           =
           \begin{mat}[r]
             1  &1  \\
             0  &-1
           \end{mat}
           \begin{mat}[r]
             3  &2  \\
             0  &1
           \end{mat}
           \begin{mat}[r]
             1  &1  \\
             0  &-1
           \end{mat}^{-1}
        \end{equation*}
      \partsitem Apakah butir sebelumnya hanya kebetulan?
        Ataukah kita selalu dapat menukar $P$ dan $P^{-1}$?
   \end{exparts}
   \begin{answer}
     \begin{exparts}
       \partsitem Pemeriksaannya mudah.
         \begin{equation*}
           \begin{mat}[r]
             1  &1  \\
             0  &-1
           \end{mat}
           \begin{mat}[r]
             3  &2  \\
             0  &1
           \end{mat}
           =
           \begin{mat}[r]
             3  &3  \\
             0  &-1
           \end{mat}
           \qquad
           \begin{mat}[r]
             3  &3  \\
             0  &-1
           \end{mat}
           \begin{mat}[r]
             1  &1  \\
             0  &-1
           \end{mat}^{-1}
           =
           \begin{mat}[r]
             3  &0  \\
             0  &1
           \end{mat}
         \end{equation*}
        \partsitem Hal itu merupakan kebetulan, dalam arti bahwa jika
          $T=PSP^{-1}$, belum tentu $T=P^{-1}SP$.
          Bahkan dalam kasus matriks diagonal~$D$, syarat $D=PTP^{-1}$ tidak
          menyiratkan bahwa $D=P^{-1}TP$.
          Matriks-matriks dari \nearbyexample{ex:DiagTwoByTwo} menunjukkannya.
          \begin{equation*}
            \begin{mat}[r]
              1  &2  \\
              1  &1
            \end{mat}
            \begin{mat}[r]
              4  &-2  \\
              1  &1
            \end{mat}
            =
            \begin{mat}[r]
              6  &0  \\
              5  &-1
            \end{mat}
            \qquad
            \begin{mat}[r]
              6  &0  \\
              5  &-1
            \end{mat}
            \begin{mat}[r]
              1  &2  \\
              1  &1
            \end{mat}^{-1}
            =
            \begin{mat}[r]
              -6  &12  \\
              -6  &11
            \end{mat}
          \end{equation*}
     \end{exparts}
   \end{answer}
  \item 
    Tunjukkan bahwa $P$ yang digunakan untuk mendiagonalkan dalam
    \nearbyexample{ex:DiagUpperTrian} tidak unik.
    \begin{answer}
      Kolom-kolom matriks tersebut adalah vektor-vektor yang terkait dengan
      nilai-nilai $x$.
      Pilihan tepatnya dan urutan pilihan tersebut bersifat sebarang.
      Sebagai contoh, kita dapat memperoleh matriks berbeda dengan menukar kedua
      kolomnya.
    \end{answer}
  \item 
    Cari rumus bagi pangkat-pangkat matriks berikut.
    \textit{Petunjuk}:~lihat \nearbyexercise{exer:PowersOfDiags}.
    \begin{equation*}
      \begin{mat}[r]
        -3  &1  \\
        -4  &2
      \end{mat}
    \end{equation*}
    \begin{answer}
      Mendiagonalkan lalu mengambil pangkat matriks diagonal menunjukkan bahwa
      \begin{equation*}
        \begin{mat}[r]
          -3  &1  \\
          -4  &2
        \end{mat}^k
        =
        \frac{1}{3}
        \begin{mat}[r]
          -1  &1  \\
          -4  &4
        \end{mat}
        +\frac{(-2)^k}{3}
        \begin{mat}[r]
           4  &-1 \\
           4  &-1
        \end{mat}.
      \end{equation*}  
     \end{answer}
  \item 
    Kita dapat bertanya bagaimana diagonalisasi berinteraksi dengan operasi
    matriks.
    Andaikan \( \map{t,s}{V}{V} \) masing-masing dapat didiagonalkan.
    Apakah \( ct \) dapat didiagonalkan bagi setiap skalar \( c \)?
    Bagaimana dengan \( t+s \)?
    \( \composed{t}{s} \)?
    \begin{answer}
      Ya, \( ct \) dapat didiagonalkan menurut teorema terakhir subbagian ini.

      Tidak, \( t+s \) belum tentu dapat didiagonalkan.
      Secara intuitif, masalah muncul ketika kedua pemetaan dapat didiagonalkan
      terhadap basis berbeda (artinya, ketika keduanya tidak
      \definend{dapat didiagonalkan secara simultan}).
      Secara khusus, kedua matriks berikut dapat didiagonalkan, tetapi jumlahnya
      tidak:
      \begin{equation*}
         \begin{mat}[r]
            1  &1  \\
            0  &0
         \end{mat}
         \qquad
         \begin{mat}[r]
           -1  &0  \\
            0  &0
         \end{mat}
      \end{equation*}
      (matriks kedua sudah diagonal; bagi matriks pertama, lihat
      \nearbyexercise{exer:DiagTheseA}).
      Jumlahnya tidak dapat didiagonalkan karena kuadratnya merupakan matriks
      nol.

     Intuisi yang sama menunjukkan bahwa \( \composed{t}{s} \) belum tentu
     dapat didiagonalkan.
     Kedua matriks berikut dapat didiagonalkan, tetapi hasil kalinya tidak:
     \begin{equation*}
        \begin{mat}[r]
           1  &0  \\
           0  &0
        \end{mat}
        \qquad
        \begin{mat}[r]
           0  &1  \\
           1  &0
        \end{mat}
     \end{equation*}
     (bagi matriks kedua, lihat \nearbyexercise{exer:DiagTheseB}).
    \end{answer}
  \item 
    Tunjukkan bahwa matriks-matriks berbentuk berikut tidak dapat didiagonalkan.
    \begin{equation*}
       \begin{mat}[r]
          1  &c  \\
          0  &1
       \end{mat}
       \qquad c\neq 0
    \end{equation*}
    \begin{answer}
      Andaikan
      \begin{equation*}
         P
         \begin{mat}
            1  &c  \\
            0  &1
         \end{mat}
         P^{-1}
         =
         \begin{mat}
            a  &0  \\
            0  &b
         \end{mat}
      \end{equation*}
      maka
      \begin{equation*}
         P
         \begin{mat}
            1  &c  \\
            0  &1
         \end{mat}
         =
         \begin{mat}
            a  &0  \\
            0  &b
         \end{mat}
         P
      \end{equation*}
      sehingga
      \begin{align*}
         \begin{mat}
            p  &q  \\
            r  &s
         \end{mat}
         \begin{mat}
            1  &c  \\
            0  &1
         \end{mat}
         &=
         \begin{mat}
            a  &0  \\
            0  &b
         \end{mat}
         \begin{mat}
            p  &q  \\
            r  &s
         \end{mat}        \\
         \begin{mat}
            p  &cp+q  \\
            r  &cr+s
         \end{mat}
         &=
         \begin{mat}
            ap  &aq  \\
            br  &bs
         \end{mat}
      \end{align*}
      Kolom pertama $P$ tidak nol karena $P$ dapat diinverskan.
      Jika \( p\neq0 \), kesamaan entri \( 1,1 \) memberikan \( a=1 \), lalu
      kesamaan entri \( 1,2 \) memberikan \( pc=0 \), yang bertentangan dengan
      \( p\neq0 \) dan \( c\neq0 \).
      Jika \( r\neq0 \), kesamaan entri \( 2,1 \) memberikan \( b=1 \), lalu
      kesamaan entri \( 2,2 \) memberikan \( rc=0 \), yang juga bertentangan
      dengan \( r\neq0 \) dan \( c\neq0 \).
      Jadi, baik \( p\neq0 \) maupun \( r\neq0 \) mustahil, sehingga kolom
      pertama $P$ nol. Ini bertentangan dengan keterbalikan $P$.
     \end{answer}
  \item 
    Tunjukkan bahwa masing-masing matriks berikut dapat didiagonalkan.
    \begin{exparts*}
      \partsitem
       \( \begin{mat}[r]
             1  &2  \\
             2  &1
          \end{mat}  \)
      \partsitem
       \( \begin{mat}
             x  &y  \\
             y  &z
           \end{mat}
           \qquad \text{$x,y,z\in\Re$}  \)
    \end{exparts*}
    \begin{answer}
      \begin{exparts}
      \partsitem Menggunakan rumus invers matriks $\nbyn{2}$ memberikan
        \begin{multline*}
           \begin{mat}
              a  &b  \\
              c  &d
           \end{mat}
           \begin{mat}[r]
              1  &2  \\
              2  &1
           \end{mat}
           \cdot\frac{1}{ad-bc}\cdot
           \begin{mat}
              d  &-b \\
             -c  &a
           \end{mat}                                  \\
           =\frac{1}{ad-bc}
           \begin{mat}
              ad+2bd-2ac-bc    &-ab-2b^2+2a^2+ab \\
              cd+2d^2-2c^2-cd  &-bc-2bd+2ac+ad
           \end{mat}
        \end{multline*}
        Sekarang pilih skalar \( a,\ldots,d \) sedemikian sehingga
        \( ad-bc\neq 0 \), \( 2d^2-2c^2=0 \), dan \( 2a^2-2b^2=0 \).
        Sebagai contoh, pilihan berikut memenuhi.
        \begin{equation*}
          \begin{mat}[r]
            1  &1  \\
            1  &-1
          \end{mat}
          \begin{mat}[r]
            1  &2  \\
            2  &1
          \end{mat}
          \cdot\frac{1}{-2}\cdot
          \begin{mat}[r]
            -1  &-1  \\
            -1  &1
          \end{mat}
          =
          \frac{1}{-2}
          \begin{mat}[r]
            -6  &0   \\
             0  &2
          \end{mat}
        \end{equation*}
      \partsitem Seperti di atas,
        \begin{multline*}
           \begin{mat}
              a  &b  \\
              c  &d
           \end{mat}
           \begin{mat}
              x  &y  \\
              y  &z
           \end{mat}
           \cdot\frac{1}{ad-bc}\cdot
           \begin{mat}
              d  &-b \\
             -c  &a
           \end{mat}                               \\
           =\frac{1}{ad-bc}
           \begin{mat}
              adx+bdy-acy-bcz    &-abx-b^2y+a^2y+abz \\
              cdx+d^2y-c^2y-cdz  &-bcx-bdy+acy+adz
           \end{mat}
        \end{multline*}
        kita mencari skalar \( a,\ldots,d \) sedemikian sehingga
        \( ad-bc\neq 0 \) dan
        \( -abx-b^2y+a^2y+abz=0 \)
        dan \( cdx+d^2y-c^2y-cdz=0 \), berapa pun nilai \( x \), \( y \), dan
        \( z \).

        Mula-mula, andaikan \( y\neq 0 \); jika tidak, matriks yang diberikan
        sudah diagonal.
        Kita akan menggunakan asumsi tersebut karena jika secara sebarang kita
        mengambil \( a=1 \), diperoleh
        \begin{align*}
           -bx-b^2y+y+bz
           &=0              \\
           (-y)b^2+(z-x)b+y
           &=0
        \end{align*}
        dan rumus kuadrat memberikan
        \begin{equation*}
           b=\frac{-(z-x)\pm\sqrt{(z-x)^2-4(-y)(y)} }{-2y}
           \qquad
           y\neq 0
        \end{equation*}
        (kedua nilai \( b \) ini real karena diskriminannya positif).
        Dengan cara yang sama, jika secara sebarang kita mengambil \( c=1 \),
        maka
        \begin{align*}
           dx+d^2y-y-dz
           &=0              \\
           (y)d^2+(x-z)d-y
           &=0
        \end{align*}
        dan diperoleh
        \begin{equation*}
           d=\frac{-(x-z)\pm\sqrt{(x-z)^2-4(y)(-y)} }{2y}
           \qquad
           y\neq 0
        \end{equation*}
        (seperti di atas, diskriminan ini positif, sehingga matriks simetris real
        \( \nbyn{2} \) serupa dengan suatu matriks diagonal real).

        Sebagai pemeriksaan, coba \( x=1 \), \( y=2 \), dan \( z=1 \).
        \begin{equation*}
           b=\frac{0\pm\sqrt{0+16} }{-4}=\mp 1
           \qquad
           d=\frac{0\pm\sqrt{0+16} }{4}=\pm 1
        \end{equation*}
        Perhatikan bahwa tidak semua dari empat pilihan
        \( (b,d)=(+1,+1),\dots,(-1,-1) \) memenuhi \( ad-bc\neq 0 \).
      \end{exparts} 
    \end{answer}
\index{dapat didiagonalkan|)}
\end{exercises}































\subsection{Nilai Eigen dan Vektor Eigen}
Selanjutnya kita akan memusatkan perhatian pada
sifat dalam \nearbylemma{lm:DiagIffBasisOfEigens}.

\begin{definition} \label{def:Eigen}
%<*df:Eigen>
Suatu transformasi \( \map{t}{V}{V} \) memiliki skalar
\definend{nilai eigen}\index{nilai eigen, vektor eigen!transformasi}%
\index{transformasi!nilai eigen, vektor eigen}
\( \lambda \)
jika terdapat \definend{vektor eigen} tak nol \( \vec{\zeta}\in V \)
sedemikian sehingga
$
  t(\vec{\zeta})=\lambda\cdot\vec{\zeta}
$.
%</df:Eigen>
\end{definition}

\noindent Dalam bahasa Jerman, ``eigen'' berarti ``ciri khas dari'' atau
``khusus bagi.''
Sebagian penulis menyebutnya nilai dan vektor \definend{karakteristik}%
\index{karakteristik!vektor, nilai}.
Tidak ada penulis yang menyebutnya vektor ``aneh.''

\begin{example}
Pemetaan proyeksi
\begin{equation*}
  \colvec{x \\ y \\ z}
     \mapsunder{\pi}
  \colvec{x \\ y \\ 0}
   \qquad x,y,z\in\C
\end{equation*}
memiliki nilai eigen \( 1 \), yang terkait dengan setiap vektor eigen
\begin{equation*}
   \colvec{x \\ y \\ 0}
\end{equation*}
di mana \( x \) dan \( y \) adalah skalar yang tidak keduanya nol.

Sebaliknya, salah satu bilangan yang bukan nilai eigen pemetaan ini adalah
\( 2 \), sebab anggapan bahwa $\pi$ menggandakan suatu vektor menghasilkan
tiga persamaan $x=2x$, $y=2y$, dan $0=2z$, sehingga tidak ada vektor
tak-$\zero$ yang digandakan.
\end{example}

Perhatikan bahwa definisi tersebut mensyaratkan vektor eigen
tak-$\zero$.
Sebagian penulis memperbolehkan $\zero$ sebagai vektor eigen untuk
$\lambda$ selama ada pula vektor tak-$\zero$ yang terkait dengan
$\lambda$.
Hal pokoknya ialah mengecualikan kasus trivial ketika $\lambda$ sedemikian
sehingga $t(\vec{v})=\lambda\vec{v}$ hanya bagi satu-satunya vektor
$\vec{v}=\zero$.

Perhatikan pula bahwa nilai eigen $\lambda$ boleh saja~$0$.
Yang menjadi persoalan ialah apakah $\vec{\zeta}$ sama dengan $\zero$.  

\begin{example}  \label{ex:NoEigenOnTrivSp}
Satu-satunya transformasi pada ruang trivial \( \set{\zero} \) adalah
%\begin{equation*}
$\zero\mapsto\zero$.
%\end{equation*}
Pemetaan ini tidak memiliki nilai eigen karena tidak ada vektor
tak-\( \zero \) $\vec{v}$ yang dipetakan ke kelipatan skalarnya sendiri,
$\lambda\cdot\vec{v}$.
\end{example}        

\begin{example} \label{ex:TransPolyOne}
Perhatikan homomorfisma \( \map{t}{\polyspace_1}{\polyspace_1} \)
yang diberikan oleh \( c_0+c_1x\mapsto(c_0+c_1)+(c_0+c_1)x \).
Meskipun kodomain $\polyspace_1$ dari $t$ berdimensi dua, daerah hasilnya
berdimensi satu, $\rangespace{t}=\set{c+cx\suchthat c\in\C}$.
Penerapan \( t \) pada suatu vektor dalam daerah hasil itu hanya akan
menskalakan ulang vektor tersebut,
\( c+cx\mapsto (2c)+(2c)x \).
Artinya, \( t \) memiliki nilai eigen \( 2 \) yang terkait dengan
vektor-vektor eigen berbentuk \( c+cx \), dengan \( c\neq 0 \).

Pemetaan ini juga memiliki nilai eigen \( 0 \) yang terkait dengan
vektor-vektor eigen berbentuk \( c-cx \), dengan \( c\neq 0 \).
\end{example}

Definisi di atas berlaku untuk pemetaan.
Kita dapat memberikan versi matriksnya.

\begin{definition}  \label{df:EigenOfMatrix}
%<*df:EigenOfMatrix>
Matriks persegi \( T \) memiliki skalar
\definend{nilai eigen}\index{nilai eigen, vektor eigen!matriks}
\( \lambda \) yang terkait dengan
\definend{vektor eigen} tak nol \( \vec{\zeta} \) jika
\( T\vec{\zeta}=\lambda\cdot\vec{\zeta} \).
%</df:EigenOfMatrix>
\end{definition}

Perluasan definisi dari pemetaan ke matriks ini wajar, tetapi ada satu hal
yang harus diperhatikan.
Nilai eigen suatu pemetaan juga merupakan nilai eigen matriks-matriks yang
merepresentasikan pemetaan tersebut, sehingga matriks-matriks yang serupa
memiliki nilai eigen yang sama.
Namun, vektor eigennya dapat berbeda\Dash matriks-matriks yang serupa tidak
harus memiliki vektor eigen yang sama.
Contoh berikut menjelaskannya.

\begin{example}
Matriks-matriks ini serupa,
\begin{equation*}
  T=
  \begin{mat}
    2  &0  \\
    0  &0
  \end{mat}
  \quad
  \hat{T}
  =
  \begin{mat}[r]
    4  &-2  \\
    4  &-2
  \end{mat}
\end{equation*}
karena $\hat{T}=PTP^{-1}$ untuk $P$ berikut.
\begin{equation*}
  P=
  \begin{mat}
    1  &1  \\
    1  &2
  \end{mat}
  \quad
  P^{-1}=
  \begin{mat}[r]
    2   &-1  \\
    -1  &1
  \end{mat}
\end{equation*}
Matriks $T$ memiliki dua nilai eigen, $\lambda_1=2$ dan~$\lambda_2=0$.
Nilai yang pertama terkait dengan vektor eigen berikut.
\begin{equation*}
  T\vec{e}_1=
  \begin{mat}
    2  &0  \\
    0  &0
  \end{mat}
  \colvec{1 \\ 0}
  =\colvec{2 \\ 0}
  =2\vec{e}_1
\end{equation*}
Misalkan $T$ merepresentasikan suatu transformasi
$\map{t}{\C^2}{\C^2}$ terhadap basis standar.
Maka aksi transformasi $t$ ini sederhana.
\begin{equation*}
  \colvec{x \\ y}\mapsunder{t}\colvec{2x \\ 0}
\end{equation*}

Tentu saja, $\hat{T}$ merepresentasikan transformasi yang sama, tetapi
terhadap basis lain, yaitu~$B$.
Kita dapat mencari basis ini.
Dengan menelusuri diagram panah dari kiri bawah ke kiri atas,
\begin{equation*}
  \begin{CD}
    V_{\wrt{\stdbasis_2}}            @>t>T>        V_{\wrt{\stdbasis_2}}       \\
    @V{\scriptstyle\identity} VV              @V{\scriptstyle\identity} VV \\
    V_{\wrt{B}}                   @>t>\hat{T}>        V_{\wrt{B}}       
  \end{CD}
\end{equation*}
terlihat bahwa $P^{-1}=\rep{\identity}{B,\stdbasis_2}$.
Menurut definisi representasi matriks suatu pemetaan, kolom pertamanya adalah
$\rep{\identity(\vec{\beta}_1)}{\stdbasis_2}=\rep{\vec{\beta}_1}{\stdbasis_2}$.
Terhadap basis standar, setiap vektor direpresentasikan oleh dirinya sendiri,
sehingga unsur basis pertama $\vec{\beta}_1$ adalah kolom pertama $P^{-1}$.
Hal yang sama berlaku bagi unsur yang lain.
\begin{equation*}
  B=\sequence{\colvec[r]{2 \\ -1},
              \colvec[r]{-1 \\ 1}
              }
\end{equation*}
Karena matriks $T$ dan~$\hat{T}$ sama-sama merepresentasikan
transformasi~$t$, keduanya mencerminkan aksi
$t(\vec{e}_1)=2\vec{e}_1$.
\begin{align*}
  &\rep{t}{\stdbasis_2,\stdbasis_2}\cdot\rep{\vec{e}_1}{\stdbasis_2}
    =T\cdot\rep{\vec{e}_1}{\stdbasis_2}                
    =2\cdot\rep{\vec{e}_1}{\stdbasis_2}                  \\               
    &\rep{t}{B,B}\cdot\rep{\vec{e}_1}{B}
    =\hat{T}\cdot\rep{\vec{e}_1}{B}                                 
    =2\cdot\rep{\vec{e}_1}{B}  
\end{align*}
Walaupun nilai eigen $2$ pada kedua persamaan tersebut sama, representasi
vektornya berbeda.
\begin{align*}
    T\cdot\rep{\vec{e}_1}{\stdbasis_2}
    =T\colvec{1 \\ 0}
    &=2\cdot\colvec{1 \\ 0}                \\
    \hat{T}\cdot\rep{\vec{e}_1}{B}  
    =\hat{T}\cdot\colvec{1 \\ 1}
    &=2\cdot\colvec{1 \\ 1}
\end{align*}
Artinya, ketika matriks yang merepresentasikan transformasi adalah
\( T=\rep{t}{\stdbasis_2,\stdbasis_2} \), matriks itu ``menganggap''
vektor-vektor kolom sebagai representasi terhadap \( \stdbasis_2 \).
Namun, \( \hat{T}=\rep{t}{B,B} \) ``menganggap'' vektor-vektor kolom
sebagai representasi terhadap \( B \), sehingga vektor-vektor kolom yang
digandakan pun berbeda.
\end{example}


% \begin{example}
% Consider again the transformation \( \map{t}{\polyspace_1}{\polyspace_1} \) 
% % from \nearbyexample{ex:TransPolyOne} 
% given by
% \( c_0+c_1x\mapsto (c_0+c_1)+(c_0+c_1)x \).
% One of its eigenvalues is \( 2 \), associated with the eigenvectors
% \( c+cx \) where \( c\neq 0 \).
% If we represent \( t \) with respect to \( B=\sequence{1+1x,1-1x} \)
% \begin{equation*}
%    T=\rep{t}{B,B}=
%    \begin{mat}[r]
%       2  &0  \\
%       0  &0
%    \end{mat}
% \end{equation*}
% then \( 2 \) is an eigenvalue of the matrix \( T \), 
% associated with these eigenvectors.
% \begin{equation*}
%    \set{\colvec{c_0 \\ c_1}\suchthat \begin{mat}[r]
%                                          2  &0  \\
%                                          0  &0
%                                       \end{mat}\colvec{c_0 \\ c_1}
%                                       =\colvec{2c_0 \\ 2c_1}  }
%   =\set{\colvec{c_0 \\ 0}\suchthat c_0\in\C,\, c_0\neq 0 }
% \end{equation*}
% On the other hand, if we represent $t$ with respect to
% \( D=\sequence{2+1x,1+0x} \)  
% \begin{equation*}
%    S=\rep{t}{D,D}=
%    \begin{mat}[r]
%       3  &1  \\
%      -3  &-1
%    \end{mat}
% \end{equation*}
% then the eigenvectors associated with the eigenvalue \( 2 \) are
% these.
% \begin{equation*}
%    \set{\colvec{c_0 \\ c_1}\suchthat \begin{mat}[r]
%                                          3  &1  \\
%                                         -3  &-1
%                                       \end{mat}\colvec{c_0 \\ c_1}
%                                       =\colvec{2c_0 \\ 2c_1}  }
%   =\set{\colvec{0 \\ c_1}\suchthat c_1\in\C,\, c_1\neq 0 }
% \end{equation*}
% \end{example}


Berikutnya kita akan melihat alat dasar untuk mencari vektor eigen dan
nilai eigen.

\begin{example}  \label{ex:IntroCharEqn}
Jika
\begin{equation*}
  T=
  \begin{mat}[r]
     1    &2    &1    \\
     2    &0    &-2   \\
    -1    &2    &3
  \end{mat}
\end{equation*}
maka, untuk mencari skalar \( x \) sedemikian sehingga
\( T\vec{\zeta}=x\vec{\zeta} \) bagi vektor eigen tak nol
\( \vec{\zeta} \), pindahkan semua suku ke ruas kiri,
\begin{equation*}
  \begin{mat}[r]
     1    &2    &1    \\
     2    &0    &-2   \\
    -1    &2    &3
  \end{mat}
  \colvec{z_1 \\ z_2 \\ z_3}
  -x\colvec{z_1 \\ z_2 \\ z_3}
  =\zero
\end{equation*}
dan faktorkan menjadi
\( (T-x I)\vec{\zeta}=\zero \).
(Perhatikan bahwa bentuknya adalah $T-xI$.
Ekspresi \( T-x \) tidak bermakna karena \( T \) adalah matriks,
sedangkan \( x \) adalah skalar.)
Sistem linear homogen ini
\begin{equation*}
  \begin{mat}
   1-x           &2            &1            \\
     2           &0-x          &-2           \\
    -1           &2            &3-x
  \end{mat}
  \colvec{z_1 \\ z_2 \\ z_3}
  =
  \colvec[r]{0 \\ 0 \\ 0}
\end{equation*}
memiliki penyelesaian tak nol $\vec{z}$ jika dan hanya jika matriksnya
singular.
Kita dapat menentukan kapan hal itu terjadi.
\begin{align*}
  0
  &=\deter{T-x I}                                               \\
  &=\begin{vmatrix}
     1-x          &2            &1            \\
     2           &0-x          &-2           \\
    -1           &2            &3-x
   \end{vmatrix}                                       \\
  &=-x^3+4x^2-4x  \\
  &=-x(x-2)^2
\end{align*}
Nilai eigennya adalah \( \lambda_1=0 \) dan \( \lambda_2=2 \).
Untuk mencari vektor eigen yang terkait, substitusikan setiap nilai eigen.
Substitusi $\lambda_1=0$ menghasilkan
\begin{equation*}
  \begin{mat}
     1-0         &2            &1            \\
     2           &0-0          &-2           \\
    -1           &2            &3-0
  \end{mat}
  \colvec{z_1 \\ z_2 \\ z_3}
  =
  \colvec[r]{0 \\ 0 \\ 0}
  \quad\Longrightarrow\quad
  \colvec{z_1 \\ z_2 \\ z_3}
  =
  \colvec[r]{a \\ -a \\ a}
\end{equation*}
untuk \( a\neq 0 \)
(\( a \) harus tak-$0$ karena vektor eigen didefinisikan sebagai
vektor tak-\( \zero \)).
Substitusi $\lambda_2=2$ menghasilkan
\begin{equation*}
  \begin{mat}
     1-2         &2            &1            \\
     2           &0-2          &-2           \\
    -1           &2            &3-2
  \end{mat}
  \colvec{z_1 \\ z_2 \\ z_3}
  =
  \colvec[r]{0 \\ 0 \\ 0}
  \quad\Longrightarrow\quad
  \colvec{z_1 \\ z_2 \\ z_3}
  =
  \colvec{b \\ 0 \\ b}
\end{equation*}
dengan \( b\neq 0 \).
\end{example}

\begin{example} \label{ex:AnotherCharPoly}
Jika
\begin{equation*}
  S=
  \begin{mat}[r]
    \pi      &1      \\
    0        &3
  \end{mat}
\end{equation*}
(di sini \( \pi \) bukan pemetaan proyeksi, melainkan bilangan
\( 3.14\ldots \)), maka
\begin{equation*}
    \begin{vmat}
      \pi-x &1         \\
      0     &3-x
    \end{vmat} 
  =
  (x-\pi)(x-3)
\end{equation*}
sehingga \( S \) memiliki nilai eigen \( \lambda_1=\pi \) dan
\( \lambda_2=3 \).
Untuk mencari vektor eigen yang terkait, mula-mula substitusikan
$\lambda_1$ untuk $x$,
\begin{equation*}
  \begin{mat}
    \pi-\pi     &1         \\
    0           &3-\pi
  \end{mat}
  \colvec{z_1 \\ z_2}
  =
  \colvec[r]{0 \\ 0}
  \qquad\Longrightarrow\qquad
  \colvec{z_1 \\ z_2}
  =
  \colvec{a \\ 0}
\end{equation*}
untuk skalar \( a\neq 0 \).
Kemudian substitusikan $\lambda_2$,
\begin{equation*}
  \begin{mat}
    \pi-3       &1         \\
    0           &3-3
  \end{mat}
  \colvec{z_1 \\ z_2}
  =
  \colvec[r]{0 \\ 0}
  \qquad\Longrightarrow\qquad
  \colvec{z_1 \\ z_2}
  =
  \colvec{-b/(\pi-3) \\ b}
\end{equation*}
di mana \( b\neq 0 \).
\end{example}

\begin{definition} \label{df:CharacteristicPoly}
%<*df:CharacteristicPoly>
\definend{Polinomial karakteristik matriks persegi}\index{karakteristik!polinomial}%
\index{matriks!polinomial karakteristik}
\( T \) adalah determinan \( \deter{T-x I} \), dengan \( x \) sebagai
variabel.
\definend{Persamaan karakteristik}\index{karakteristik!persamaan}%
\index{matriks!polinomial karakteristik}
adalah $\deter{T-xI}=0$.
\definend{Polinomial karakteristik transformasi} \( t \) adalah
polinomial karakteristik dari sembarang representasi matriks
\( \rep{t}{B,B} \).\index{transformasi!polinomial karakteristik}
%</df:CharacteristicPoly>
\end{definition}

\noindent 
%<*df:CharacteristicPolyExer>
Polinomial karakteristik matriks \( \nbyn{n} \), atau transformasi
\( \map{t}{\C^n}{\C^n} \), berderajat~\( n \).
\nearbyexercise{exer:CharPolyTransWellDefed} memeriksa bahwa polinomial
karakteristik suatu transformasi terdefinisi dengan baik; artinya,
polinomial karakteristiknya sama untuk basis representasi mana pun yang
kita gunakan.
%</df:CharacteristicPolyExer>

\begin{lemma} \label{le:MapNonTrivSpHasEigen}
%<*lm:MapNonTrivSpHasEigen>
Transformasi linear pada ruang vektor nontrivial memiliki sedikitnya satu
nilai eigen.
%</lm:MapNonTrivSpHasEigen>
\end{lemma}

\begin{proof}
%<*pf:MapNonTrivSpHasEigen>
Setiap akar polinomial karakteristik merupakan nilai eigen.
Di atas bilangan kompleks, setiap polinomial berderajat satu atau lebih
memiliki akar.
%</pf:MapNonTrivSpHasEigen>
\end{proof}

\begin{remark}
Hasil tersebut merupakan alasan kita menggunakan skalar bilangan kompleks
dalam bab ini.
Seandainya kita tetap menggunakan skalar bilangan real, akan ada polinomial
karakteristik, seperti $x^2+1$, yang tidak dapat difaktorkan.
\end{remark}

% Notice the familiar form of the sets of eigenvectors in the above examples.

\begin{definition} \label{df:Eigenspace}
%<*df:Eigenspace>
\definend{Ruang eigen transformasi~$t$ yang terkait dengan nilai
eigen~$\lambda$}\index{ruang eigen}\index{transformasi!ruang eigen}
adalah
$
  V_{\lambda}=\set{\vec{\zeta}\suchthat t(\vec{\zeta}\,)=\lambda\vec{\zeta}\,}
              % \union\set{\zero\,}
$.
Ruang eigen suatu matriks didefinisikan secara analog.
%</df:Eigenspace>
\end{definition}

% \begin{remark}
% In \nearbydefinition{def:Eigen} we specified that eigenvectors must be
% non-$\zero$.
% So, strictly speaking, the eigenspace for $\lambda$ contains the set of
% associated eigenvectors along with $\zero$.
% We need the zero vector for the next result. 
% \end{remark}

\begin{lemma}  \label{le:EigSpaceIsSubSp}
%<*lm:EigSpaceIsSubSp>
Ruang eigen adalah subruang.
Ruang eigen merupakan subruang nontrivial.
%</lm:EigSpaceIsSubSp>
\end{lemma}

\begin{proof}
%<*pf:EigSpaceIsSubSp>
Pertama, perhatikan bahwa $V_{\lambda}$ tidak kosong; ruang ini memuat
vektor nol karena $t(\zero)=\zero$, yang sama dengan
$\lambda\cdot \zero$.
Untuk menunjukkan bahwa ruang eigen merupakan subruang, tinggal diperiksa
bahwa himpunan ini tertutup terhadap kombinasi linear.
Ambil \( \vec{\zeta}_1,\ldots,\vec{\zeta}_n\in V_{\lambda} \). Maka
\begin{align*}
  t(\lincombo{c}{\vec{\zeta}})
  &=c_1t(\vec{\zeta}_1)+\dots+c_nt(\vec{\zeta}_n)               \\
  &=c_1\lambda\vec{\zeta}_1+\dots+c_n\lambda\vec{\zeta}_n          \\
  &=\lambda(c_1\vec{\zeta}_1+\dots+c_n\vec{\zeta}_n)
\end{align*}
menunjukkan bahwa kombinasi tersebut juga merupakan anggota $V_{\lambda}$.
% (despite that the zero vector isn't an eigenvector, 
% the second equality holds even if some \( \vec{\zeta}_i \) is \( \zero \) since
% \( t(\zero)=\lambda\cdot\zero=\zero \)).

Ruang~$V_\lambda$ memuat lebih dari sekadar vektor nol karena, menurut
definisi, $\lambda$ merupakan nilai eigen hanya jika
$t(\vec{\zeta}\,)=\lambda\vec{\zeta}$ memiliki penyelesaian
$\vec{\zeta}$ selain $\zero$.
%</pf:EigSpaceIsSubSp>
\end{proof}

\begin{example}   \label{ex:AlgMultDifferentThanGeoMult}
Berikut adalah ruang-ruang eigen yang terkait dengan nilai eigen \( 0 \)
dan \( 2 \) dalam \nearbyexample{ex:IntroCharEqn}.
\begin{equation*}
  V_0=\set{\colvec[r]{a \\ -a \\ a}\suchthat a\in\C},
  \qquad
  V_2=\set{\colvec{b \\ 0 \\ b}\suchthat b\in\C}.
\end{equation*}
\end{example}

\begin{example} \label{ex:AnotherCharPolyReprise}
Berikut adalah ruang-ruang eigen untuk nilai eigen \( \pi \) dan~\( 3 \)
dalam \nearbyexample{ex:AnotherCharPoly}.
\begin{equation*}
  V_{\pi}=\set{\colvec{a \\ 0}\suchthat a\in\C}
  \qquad
  V_3=\set{\colvec{-b/(\pi-3) \\ b}\suchthat b\in\C}
\end{equation*}
\end{example}

Persamaan karakteristik dalam \nearbyexample{ex:IntroCharEqn} adalah
\( 0=x(x-2)^2 \), sehingga dalam arti tertentu \( 2 \) merupakan nilai
eigen sebanyak dua kali.
Namun, jumlah vektor eigennya bukan dua kali lipat: dimensi ruang eigen
$V_2$ yang terkait adalah satu, bukan dua.
Contoh berikut menunjukkan kasus ketika suatu bilangan merupakan akar
ganda persamaan karakteristik dan dimensi ruang eigen yang terkait adalah
dua.

\begin{example}  \label{ex:EigenvaluesProjection}
Terhadap basis-basis standar, matriks berikut
\begin{equation*}
  \begin{mat}[r]
     1  &0  &0  \\
     0  &1  &0  \\
     0  &0  &0
  \end{mat}
\end{equation*}
merepresentasikan proyeksi.
\begin{equation*}
  \colvec{x \\ y \\ z}
     \mapsunder{\pi}
  \colvec{x \\ y \\ 0}
   \qquad x,y,z\in\C
\end{equation*}
Persamaan karakteristiknya
\begin{align*}
  0
  &=\deter{T-x I}                                     \\
  &=\begin{vmatrix}
     1-x          &0            &0            \\
     0           &1-x          &0           \\
     0           &0            &0-x
   \end{vmatrix}                                       \\
  &=(1-x)^2(0-x)
\end{align*}
memiliki akar ganda~$x=1$ beserta akar tunggal~$x=0$.
Ruang eigen yang terkait dengan nilai eigen \( 1 \) dan ruang eigen yang
terkait dengan nilai eigen \( 0 \) mudah dicari.
\begin{equation*}
   V_1=\set{\colvec{c_1 \\ c_2 \\ 0}\suchthat c_1,c_2\in\C}
   \qquad
   V_0=\set{\colvec{0 \\ 0 \\ c_3}\suchthat c_3\in\C}
\end{equation*}
Perhatikan bahwa $V_1$ berdimensi dua.
\end{example}

\begin{definition}
Jika suatu polinomial karakteristik difaktorkan menjadi
$(x-\lambda_1)^{m_1}\cdots (x-\lambda_k)^{m_k}$, maka nilai eigen
$\lambda_i$ memiliki \definend{multiplisitas aljabar}~$m_i$%
\index{multiplisitas!aljabar}\index{multiplisitas aljabar}.
\definend{Multiplisitas geometris}\index{multiplisitas!geometris}%
\index{multiplisitas geometris}
nilai eigen itu adalah dimensi ruang eigen~$V_{\lambda_i}$ yang terkait.  
\end{definition}

Dalam \nearbyexample{ex:EigenvaluesProjection} terdapat dua nilai eigen.
Untuk $\lambda_1=1$, multiplisitas aljabar dan geometrisnya sama-sama~$2$.
Untuk $\lambda_2=0$, kedua multiplisitas itu sama-sama~$1$.
Sebaliknya, \nearbyexample{ex:AlgMultDifferentThanGeoMult} menunjukkan
bahwa nilai eigen $\lambda=2$ memiliki multiplisitas aljabar~$2$, tetapi
multiplisitas geometris~$1$.
Untuk setiap transformasi, tiap nilai eigen memiliki multiplisitas
geometris yang lebih besar daripada atau sama dengan~$1$, berdasarkan
\nearbylemma{le:EigSpaceIsSubSp}.
(Selain itu, multiplisitas geometris suatu nilai eigen harus lebih kecil
daripada atau sama dengan multiplisitas aljabarnya, walaupun pembuktiannya
berada di luar cakupan kita.)

Berdasarkan \nearbylemma{le:EigSpaceIsSubSp}, jika dua vektor eigen
$\vec{v}_1$ dan $\vec{v}_2$ terkait dengan nilai eigen yang sama, maka
kombinasi linear keduanya juga merupakan vektor eigen yang terkait dengan
nilai eigen tersebut.
Sebagai ilustrasi, dengan merujuk pada contoh sebelumnya, jumlah dua
anggota $V_1$ berikut
\begin{equation*}
  \colvec[r]{1 \\ 0 \\ 0}+\colvec[r]{0 \\ 1 \\ 0}
\end{equation*}
menghasilkan anggota lain dari $V_1$.

Hasil berikut membahas keadaan ketika vektor-vektornya berasal dari
ruang-ruang eigen yang berbeda.

\begin{theorem}  \label{th:DistEValueGivesLIEvecs}
%<*th:DistEValueGivesLIEvecs>
Untuk sembarang himpunan nilai eigen yang berbeda-beda dari suatu pemetaan
atau matriks, himpunan vektor eigen terkait yang memuat satu vektor untuk
setiap nilai eigen adalah bebas linear.
%</th:DistEValueGivesLIEvecs>
\end{theorem}

\begin{proof}
%<*pf:DistEValueGivesLIEvecs0>
Kita akan menggunakan induksi pada banyaknya nilai eigen.
Pada langkah dasar, banyaknya nilai eigen adalah nol.
Dengan demikian, himpunan vektor yang terkait kosong dan karenanya bebas
linear.
% If there is 
% only one eigenvalue then the set of associated eigenvectors is
% a singleton set with a non-$\zero$ member, 
% and so is linearly independent.
%</pf:DistEValueGivesLIEvecs0>

%<*pf:DistEValueGivesLIEvecs1>
Untuk langkah induksi, anggap pernyataan tersebut benar bagi setiap
himpunan yang terdiri atas \( k\geq 0 \) nilai eigen berbeda.
Perhatikan nilai-nilai eigen berbeda
\( \lambda_1,\dots,\lambda_{k+1} \), dan misalkan
\( \vec{v}_1,\dots,\vec{v}_{k+1} \) adalah vektor-vektor eigen yang
terkait.
Andaikan
$\zero=c_1\vec{v}_1+\dots+c_k\vec{v}_k+c_{k+1}\vec{v}_{k+1}$. 
Dari persamaan itu, turunkan dua persamaan. Persamaan pertama diperoleh
dengan mengalikan kedua ruas dengan \( \lambda_{k+1} \),
$\zero=c_1\lambda_{k+1}\vec{v}_1+\dots+c_{k+1}\lambda_{k+1}\vec{v}_{k+1}$
dan persamaan kedua dengan menerapkan pemetaan pada kedua ruas,
$\zero=c_1t(\vec{v}_1)+\dots+c_{k+1}t(\vec{v}_{k+1})
  =c_1\lambda_1\vec{v}_1+\dots+c_{k+1}\lambda_{k+1}\vec{v}_{k+1}$
(penerapan matriks memberikan hasil yang sama).
Kurangkan persamaan kedua dari persamaan pertama.
\begin{equation*}
  \zero=
  c_1(\lambda_{k+1}-\lambda_1)\vec{v}_1+\dots
  +c_k(\lambda_{k+1}-\lambda_k)\vec{v}_k
      +c_{k+1}(\lambda_{k+1}-\lambda_{k+1})\vec{v}_{k+1}
\end{equation*}
Suku $\vec{v}_{k+1}$ lenyap.
Hipotesis induksi kemudian memberikan
\( c_1(\lambda_{k+1}-\lambda_1)=0 \), \ldots, \( c_k(\lambda_{k+1}-\lambda_k)=0 \).
Karena nilai-nilai eigennya berbeda, semua koefisien
\( c_1,\,\dots,\,c_k \) adalah~\( 0 \).
Dengan demikian, yang tersisa adalah persamaan
\( \zero=c_{k+1}\vec{v}_{k+1} \), sehingga \( c_{k+1} \) juga~\( 0 \).
%</pf:DistEValueGivesLIEvecs1>
\end{proof}

\begin{example}
Nilai-nilai eigen matriks
\begin{equation*}
     \begin{mat}[r]
        2   &-2   &2   \\
        0   &1    &1   \\
       -4   &8    &3
     \end{mat}
\end{equation*}
berbeda-beda: \( \lambda_1=1 \), \( \lambda_2=2 \), dan~\( \lambda_3=3 \).
Salah satu himpunan vektor eigen yang terkait,
\begin{equation*}
  \set{
       \colvec[r]{2 \\ 1 \\ 0},
       \colvec[r]{9 \\ 4 \\ 4},
       \colvec[r]{2 \\ 1 \\ 2}  }
\end{equation*}
adalah bebas linear.
\end{example}

\begin{corollary} \label{co:DistinctEivenvaluesImpliesDiagonal}
%<*co:DistinctEivenvaluesImpliesDiagonal>
Matriks \( \nbyn{n} \) yang memiliki \( n \) nilai eigen berbeda dapat
didiagonalkan.
%</co:DistinctEivenvaluesImpliesDiagonal>
\end{corollary}

\begin{proof}
%<*pf:DistinctEivenvaluesImpliesDiagonal>
Bentuklah basis yang terdiri atas vektor-vektor eigen.
Terapkan \nearbylemma{lm:DiagIffBasisOfEigens}.
%</pf:DistinctEivenvaluesImpliesDiagonal>
\end{proof}

\begin{example} \label{ex:SummaryOfDiagonalizable}
Dalam contoh sebelumnya, kita menunjukkan bahwa matriks
\begin{equation*}
     T=\begin{mat}[r]
        2   &-2   &2   \\
        0   &1    &1   \\
       -4   &8    &3
     \end{mat}
\end{equation*}
dapat didiagonalkan dengan nilai eigen
$\lambda_1=1$, $\lambda_2=2$, dan~$\lambda_3=3$.
Berikut adalah vektor-vektor eigen terkait, yang membentuk basis~$B$.
\begin{equation*}
       \vec{\beta}_1=\colvec[r]{2 \\ 1 \\ 0}\quad
       \vec{\beta}_2=\colvec[r]{9 \\ 4 \\ 4}\quad
       \vec{\beta}_3=\colvec[r]{2 \\ 1 \\ 2} 
\end{equation*}
Diagram panah
\begin{equation*}
  \begin{CD}
    V_{\wrt{\stdbasis_3}}            @>t>T>        V_{\wrt{\stdbasis_3}}       \\
    @V{\scriptstyle\identity} VV              @V{\scriptstyle\identity} VV \\
    V_{\wrt{B}}                   @>t>D>        V_{\wrt{B}}       
  \end{CD}
\end{equation*}
memberikan persamaan berikut, dengan
$P=\rep{\identity}{\stdbasis_3,B}$.
\begin{align*}
  D &= PTP^{-1}  \\
  \begin{mat}
    1  &0  &0 \\
    0  &2  &0 \\
    0  &0  &3
  \end{mat}
  &=
    \begin{mat}
      -2  &5  &-1/2  \\
       1  &-2  &0    \\
      -2  &4  &1/2 
    \end{mat}
     \begin{mat}[r]
        2   &-2   &2   \\
        0   &1    &1   \\
       -4   &8    &3
     \end{mat}
    \begin{mat}
      2  &9  &2  \\
      1  &4  &1  \\
      0  &4  &2
    \end{mat}
\end{align*}
% sage: T = matrix(QQ, [[2,-2,2], [0,1,1], [-4,8,3]])
% sage: T
% [ 2 -2  2]
% [ 0  1  1]
% [-4  8  3]
% sage: T.eigenvectors_right()
% [(3, [
%   (1, 1/2, 1)
%   ], 1), (2, [
%   (1, 4/9, 4/9)
%   ], 1), (1, [
%   (1, 1/2, 0)
%   ], 1)]
% sage: P_inv = matrix(QQ, [[2,9,2], [1,4,1], [0,4,2]])
% sage: P_inv.inverse()
% [  -2    5 -1/2]
% [   1   -2    0]
% [  -2    4  1/2]
Pada bagian bawah diagram berlaku
$t(\vec{\beta}_1)=1\cdot\vec{\beta}_1$,
$t(\vec{\beta}_2)=2\cdot\vec{\beta}_2$, dan
$t(\vec{\beta}_3)=3\cdot\vec{\beta}_3$.
Dengan menulis ulang, pemetaan $t-1$ mengirim $\vec{\beta}_1$
ke~$\zero$, pemetaan $t-2$ mengirim $\vec{\beta}_2$ ke~$\zero$, dan
$(t-3)(\vec{\beta}_3)=\zero$.
Artinya, aksinya pada $B$ adalah sebagai berikut.
\begin{equation*}
  \vec{\beta}_1\mapsunder{t-1}\zero  \qquad
  \vec{\beta}_2\mapsunder{t-2}\zero  \qquad
  \vec{\beta}_3\mapsunder{t-3}\zero 
\end{equation*}
Beralih ke representasi dalam diagram panah tersebut, baris atas menyatakan
bahwa matriks $T-(1\cdot I)$, yang merepresentasikan $t-1$ terhadap
$\stdbasis_3,\stdbasis_3$, jika dikalikan dengan representasi
$\vec{\beta}_1$ terhadap~$\stdbasis_3$, akan menghasilkan nol.
\begin{equation*}
  \begin{mat}
    1  &-2  &2  \\
    0  &0   &1  \\
   -4  &8   &2  \\
  \end{mat}
  \colvec{2 \\ 1 \\ 0}
  =
  \colvec{0  \\ 0 \\ 0}
\end{equation*}
Serupa dengan itu, matriks $T-(2\cdot I)$, yang merepresentasikan $t-2$
terhadap $\stdbasis_3,\stdbasis_3$, menghasilkan nol ketika dikalikan
dengan representasi $\vec{\beta}_2$ terhadap~$\stdbasis_3$.
\begin{equation*}
  \begin{mat}
    0  &-2  &2  \\
    0  &-1   &1  \\
   -4  &8   &1  \\
  \end{mat}
  \colvec{9 \\ 4 \\ 4}
  =
  \colvec{0  \\ 0 \\ 0}
\end{equation*}
Dan tentu saja, matriks $T-(3\cdot I)$ yang dikalikan dengan representasi
$\vec{\beta}_3$ terhadap~$\stdbasis_3$ juga menghasilkan nol.
\end{example}

Singkatnya, bagian ini menunjukkan bahwa sebagian matriks serupa dengan
matriks diagonal.
Gagasan nilai eigen muncul sebagai entri-entri matriks diagonal tersebut,
walaupun definisinya berlaku lebih luas, tidak hanya bagi matriks yang
dapat didiagonalkan.
Untuk mencari nilai eigen, kita mendefinisikan persamaan karakteristik,
yang kemudian mengantarkan kita pada hasil terakhir: suatu kriteria agar
matriks dapat didiagonalkan.
(Walaupun berguna bagi teori, perhatikan bahwa dalam penerapan, mencari
nilai eigen dengan cara ini biasanya tidak praktis; salah satu alasannya,
matriksnya mungkin besar dan akar polinomial berderajat tinggi sulit
dicari.)

Pada bagian berikutnya, kita mempelajari matriks yang tidak dapat
didiagonalkan.


\begin{exercises}
  \item Matriks ini mempunyai dua nilai eigen $\lambda_1=3$, $\lambda_2=-4$.
    \begin{equation*}
      \begin{mat}
        4  &1 \\
        -8 &-5
      \end{mat}
    \end{equation*}
    Berikan dua matriks berbentuk diagonal berbeda yang serupa dengannya.
    \begin{answer}
      Kita dapat menukar urutan kolom dengan menukar urutan basis yang
      digunakan untuk membuat representasi.
      \begin{equation*}
        \begin{mat} 
          3 &0 \\ 
          0 &-4
        \end{mat}
        \qquad
        \begin{mat} 
          -4 &0 \\ 
          0 &3
        \end{mat}
      \end{equation*}
    \end{answer}
  \item 
    Untuk setiap matriks, tentukan polinomial karakteristik dan nilai eigennya.
    \begin{exparts*}
      \partsitem \( \begin{mat}[r]
                 10  &-9 \\
                  4  &-2
            \end{mat}  \)
      \partsitem $\begin{mat}[r]
                    1  &2  \\
                    4  &3  
                 \end{mat}$
      \partsitem \( \begin{mat}[r]
                  0  &3  \\
                  7 &0
                 \end{mat} \)
      \partsitem \( \begin{mat}[r]  
                  0  &0  \\
                  0  &0
            \end{mat}  \)
      \partsitem \( \begin{mat}[r]
                  1  &0  \\
                  0  &1
            \end{mat}  \)
    \end{exparts*}
    \begin{answer}
      \begin{exparts}
         \partsitem Persamaan
           \begin{equation*}
             0=
             \begin{vmatrix}
               10-x  &-9  \\
               4     &-2-x
             \end{vmatrix}
             =(10-x)(-2-x)-(-36)
           \end{equation*}
           menyederhana menjadi persamaan karakteristik \( x^2-8x+16=0 \). 
           Karena persamaan itu terfaktorkan menjadi $(x-4)^2$, hanya ada
           satu nilai eigen, yaitu \( \lambda_1=4 \).
         \partsitem $0=(1-x)(3-x)-8=x^2-4x-5$; $\lambda_1=5$, $\lambda_2=-1$
         \partsitem \( x^2-21=0 \); 
           \( \lambda_1=\sqrt{21} \), $\lambda_2=-\sqrt{21}$
         \partsitem \( x^2=0 \); \( \lambda_1=0 \)
         \partsitem \( x^2-2x+1=0 \); \( \lambda_1=1 \)
       \end{exparts}  
     \end{answer}
  \recommended \item
    Untuk setiap matriks, tentukan persamaan karakteristik, nilai eigen,
    dan vektor eigen yang terkait.
    \begin{exparts*}
      \partsitem \( \begin{mat}[r]
                  3  &0  \\
                  8  &-1
            \end{mat}  \)
      \partsitem \( \begin{mat}[r]
                  3  &2  \\
                 -1  &0
            \end{mat}  \)
    \end{exparts*}
    \begin{answer}
       \begin{exparts}
         \partsitem Persamaan karakteristiknya adalah \( (3-x)(-1-x)=0 \).
           Akar-akarnya, yaitu nilai eigennya, adalah \( \lambda_1=3 \) dan 
           \( \lambda_2=-1 \).
           Untuk mencari vektor eigen, kita tinjau persamaan ini.
           \begin{equation*}
             \begin{mat}
               3-x  &0    \\
               8    &-1-x
             \end{mat}
             \colvec{b_1  \\  b_2}
             =\colvec[r]{0  \\  0}
           \end{equation*}
           Untuk vektor eigen yang terkait dengan $\lambda_1=3$,
           kita tinjau sistem linear yang dihasilkan.
           \begin{equation*}
             \begin{linsys}{2}
               0\cdot b_1  &+  &0\cdot b_2  &=  &0  \\
               8\cdot b_1  &+  &-4\cdot b_2 &=  &0
             \end{linsys}
           \end{equation*}
           Ruang eigennya adalah himpunan vektor yang komponen keduanya 
           dua kali komponen pertama.
           \begin{equation*}
             \set{\colvec{b_2/2 \\ b_2}\suchthat b_2\in\C}
             \qquad
             \begin{mat}[r]
               3  &0  \\
               8  &-1
             \end{mat}
             \colvec{b_2/2  \\ b_2}
             =3\cdot\colvec{b_2/2 \\ b_2}
           \end{equation*}
           (Di sini, parameternya adalah $b_2$ semata-mata karena itulah variabel
           bebas dalam sistem di atas.)
           Jadi, berikut ini adalah vektor eigen yang terkait dengan nilai eigen $3$.
           \begin{equation*}
             \colvec[r]{1 \\ 2}
           \end{equation*}

           Vektor eigen yang terkait dengan $\lambda_2=-1$ dicari dengan cara serupa.
           Sistem ini
           \begin{equation*}
             \begin{linsys}{2}
               4\cdot b_1  &+  &0\cdot b_2  &=  &0  \\
               8\cdot b_1  &+  &0\cdot b_2 &=  &0
             \end{linsys}
           \end{equation*}
           menghasilkan himpunan vektor yang komponen pertamanya 
           nol.
           \begin{equation*}
             \set{\colvec{0 \\ b_2}\suchthat b_2\in\C}
             \qquad
             \begin{mat}[r]
               3  &0  \\
               8  &-1
             \end{mat}
             \colvec{0  \\ b_2}
             =-1\cdot\colvec{0 \\ b_2}
           \end{equation*}
           Jadi, berikut ini adalah vektor eigen yang terkait dengan $\lambda_2$.
           \begin{equation*}
             \colvec[r]{0 \\ 1}
           \end{equation*}
         \partsitem Persamaan karakteristiknya adalah
           \begin{equation*}
             0=
             \begin{vmatrix}
               3-x  &2  \\
               -1   &-x               
             \end{vmatrix}
             =x^2-3x+2=(x-2)(x-1)
           \end{equation*}
           sehingga nilai eigennya adalah $\lambda_1=2$ dan $\lambda_2=1$.
           Untuk mencari vektor eigen, tinjau sistem ini.
           \begin{equation*}
             \begin{linsys}{2}
               (3-x)\cdot b_1  &+  &2\cdot b_2  &=  &0  \\
               -1\cdot b_1     &-  &x\cdot b_2  &=  &0
             \end{linsys}
           \end{equation*}
           Untuk $\lambda_1=2$ kita memperoleh 
           \begin{equation*}
             \begin{linsys}{2}
                1\cdot b_1  &+  &2\cdot b_2  &=  &0  \\
               -1\cdot b_1  &-  &2\cdot b_2  &=  &0
             \end{linsys}
           \end{equation*}
           yang menghasilkan ruang eigen dan vektor eigen berikut.
           \begin{equation*}
             \set{\colvec{-2b_2 \\ b_2}
                   \suchthat b_2\in\C}
             \qquad
             \colvec[r]{-2 \\ 1}
           \end{equation*}
           Untuk $\lambda_2=1$, sistemnya adalah 
           \begin{equation*}
             \begin{linsys}{2}
                2\cdot b_1  &+  &2\cdot b_2  &=  &0  \\
               -1\cdot b_1  &-  &1\cdot b_2  &=  &0
             \end{linsys}
           \end{equation*}
           yang menghasilkan berikut ini.
           \begin{equation*}
             \set{\colvec{-b_2 \\ b_2}
                   \suchthat b_2\in\C}
             \qquad
             \colvec[r]{-1 \\ 1}
           \end{equation*}
       \end{exparts}  
     \end{answer}
  \item
    Tentukan persamaan karakteristik, nilai eigen, dan vektor eigen yang
    terkait untuk matriks ini.
    \textit{Petunjuk.}
      Nilai eigennya kompleks.
    \begin{equation*}
      \begin{mat}[r]
         -2  &-1 \\
          5  &2
      \end{mat}  
    \end{equation*}
    \begin{answer}
         Persamaan karakteristik 
           \begin{equation*}
             0=
             \begin{vmatrix}
               -2-x  &-1  \\
               5     &2-x               
             \end{vmatrix}
             =x^2+1
           \end{equation*}
           mempunyai akar kompleks $\lambda_1=i$ dan $\lambda_2=-i$.
           Tinjau sistem berikut. 
           \begin{equation*}
             \begin{linsys}{2}
               (-2-x)\cdot b_1  &-  &1\cdot b_2      &=  &0  \\
               5\cdot b_1       &   &(2-x)\cdot b_2  &=  &0
             \end{linsys}
           \end{equation*}
           Untuk $\lambda_1=i$, Metode Gauss memberikan reduksi berikut.
           \begin{equation*}
             \begin{linsys}{2}
                (-2-i)\cdot b_1  &-  &1\cdot b_2      &=  &0  \\
                 5\cdot b_1       &+  &(2-i)\cdot b_2  &=  &0
             \end{linsys}
             \grstep{(-5/(-2-i))\rho_1+\rho_2}
             \begin{linsys}{2}
                (-2-i)\cdot b_1  &-  &1\cdot b_2      &=  &0  \\
                                 &   &0               &=  &0
             \end{linsys}
           \end{equation*}
           (Untuk perhitungan di kanan bawah, samakan penyebut
           \begin{equation*}
             \frac{5}{-2-i}+(2-i)
             =
             \frac{5}{-2-i}+\frac{-2-i}{-2-i}\cdot (2-i)
             =
             \frac{5+(-5)}{-2-i}
           \end{equation*}
           untuk melihat bahwa hasilnya adalah persamaan $0=0$.)
           Berikut ruang eigen dan vektor eigen yang dihasilkan.
           \begin{equation*}
             \set{\colvec{(1/(-2-i))b_2 \\ b_2}
                   \suchthat b_2\in\C}
             \qquad
             \colvec{1/(-2-i) \\ 1}
           \end{equation*}
           Untuk $\lambda_2=-i$, sistem 
           \begin{equation*}
             \begin{linsys}{2}
                (-2+i)\cdot b_1  &-  &1\cdot b_2      &=  &0  \\
                 5\cdot b_1       &+  &(2+i)\cdot b_2  &=  &0
             \end{linsys}
             \grstep{(-5/(-2+i))\rho_1+\rho_2}
             \begin{linsys}{2}
                (-2+i)\cdot b_1  &-  &1\cdot b_2      &=  &0  \\
                                 &   &0               &=  &0
             \end{linsys}
           \end{equation*}
           menghasilkan berikut ini.
           \begin{equation*}
             \set{\colvec{(1/(-2+i))b_2 \\ b_2}
                   \suchthat b_2\in\C}
             \qquad
             \colvec{1/(-2+i) \\ 1}
           \end{equation*}
    \end{answer}
  \item  
    Tentukan polinomial karakteristik, nilai eigen, dan vektor eigen yang
    terkait untuk matriks ini.
    \begin{equation*}
      \begin{mat}[r]
        1  &1  &1  \\
        0  &0  &1  \\
        0  &0  &1
      \end{mat}
    \end{equation*}
    \begin{answer}
      Persamaan karakteristiknya adalah
      \begin{equation*}
        0=
        \begin{vmatrix}
          1-x  &1   &1   \\
          0    &-x  &1   \\
          0    &0   &1-x
        \end{vmatrix}
        =(1-x)^2(-x)
      \end{equation*}
      sehingga nilai eigennya adalah $\lambda_1=1$ (ini merupakan akar berulang
      dari persamaan tersebut) dan $\lambda_2=0$.
      Selanjutnya, tinjau sistem ini.
      \begin{equation*}
        \begin{linsys}{3}
          (1-x)\cdot b_1  &+  &b_2         &+  &b_3            &=  &0  \\
                          &   &-x\cdot b_2 &+  &b_3            &=  &0  \\
                          &   &            &   &(1-x)\cdot b_3 &= &0  
        \end{linsys}
      \end{equation*}
      Ketika $x=\lambda_1=1$, himpunan penyelesaiannya adalah ruang eigen berikut.
      \begin{equation*}
        \set{\colvec{b_1 \\ 0 \\ 0}\suchthat b_1\in\C}
      \end{equation*}
      Ketika $x=\lambda_2=0$, himpunan penyelesaiannya adalah ruang eigen berikut.
      \begin{equation*}
        \set{\colvec{-b_2 \\ b_2 \\ 0}\suchthat b_2\in\C}
      \end{equation*}
      Jadi, berikut adalah vektor eigen yang terkait dengan $\lambda_1=1$ dan 
      $\lambda_2=0$.
      \begin{equation*}
        \colvec[r]{1 \\ 0 \\ 0}
        \qquad
        \colvec[r]{-1 \\ 1 \\ 0}  
      \end{equation*}
    \end{answer}
  \recommended \item
    Untuk setiap matriks, tentukan persamaan karakteristik, nilai eigen,
    dan vektor eigen yang terkait.
    \begin{exparts*}
      \partsitem \( \begin{mat}[r]
              3  &-2 &0  \\
             -2  &3  &0  \\
              0  &0  &5
            \end{mat}  \)
      \partsitem \( \begin{mat}[r]
              0  &1   &0  \\
              0  &0   &1  \\
              4  &-17 &8
            \end{mat}  \)
    \end{exparts*}
    \begin{answer}
      \begin{exparts}
         \partsitem Persamaan karakteristiknya adalah
           \begin{equation*}
             0=
             \begin{vmatrix}
              3-x  &-2   &0  \\
             -2    &3-x  &0  \\
              0    &0    &5-x
             \end{vmatrix}
             =-x^3+11x^2-35x+25=-1\cdot(x-1)(x-5)^2
           \end{equation*}
           sehingga nilai eigennya adalah $\lambda_1=1$ dan nilai eigen
           berulang $\lambda_2=5$.
           Untuk mencari vektor eigen, tinjau sistem ini.
           \begin{equation*}
             \begin{linsys}{3}
               (3-x)\cdot b_1  &-  &2\cdot b_2      &   &   &=  &0  \\
               -2\cdot b_1     &+  &(3-x)\cdot b_2  &   &   &=  &0  \\
                               &   &                &   &(5-x)\cdot b_3 &= &0 
             \end{linsys}
           \end{equation*}
           Untuk $\lambda_1=1$ kita memperoleh 
           \begin{equation*}
             \begin{linsys}{3}
                2\cdot b_1     &-  &2\cdot b_2   &   &   &=  &0  \\
               -2\cdot b_1     &+  &2\cdot b_2   &   &   &=  &0  \\
                               &   &             &   &4\cdot b_3 &= &0 
             \end{linsys}
           \end{equation*}
           yang menghasilkan ruang eigen dan vektor eigen berikut.
           \begin{equation*}
             \set{\colvec{b_2 \\ b_2 \\ 0}
                   \suchthat b_2\in\C}
             \qquad
             \colvec[r]{1 \\ 1 \\ 0}
           \end{equation*}
           Untuk $\lambda_2=5$, sistemnya adalah 
           \begin{equation*}
             \begin{linsys}{3}
                -2\cdot b_1  &-  &2\cdot b_2   &   &   &=  &0  \\
               -2\cdot b_1   &-  &2\cdot b_2   &   &   &=  &0  \\
                             &   &             &   &0\cdot b_3 &= &0 
             \end{linsys}
           \end{equation*}
           yang menghasilkan berikut ini.
           \begin{equation*}
             \set{\colvec{-b_2 \\ b_2 \\ 0}+\colvec{0 \\ 0 \\ b_3}
                   \suchthat b_2,b_3\in\C}
             \qquad
             \colvec[r]{-1 \\ 1 \\ 0},\,\colvec[r]{0 \\ 0 \\ 1}
           \end{equation*}
         \partsitem Persamaan karakteristiknya adalah
           \begin{equation*}
             0=
             \begin{vmatrix}
              -x   &1    &0  \\
              0    &-x   &1  \\
              4    &-17  &8-x
             \end{vmatrix}
             =-x^3+8x^2-17x+4=-1\cdot(x-4)(x^2-4x+1)
           \end{equation*}
           dan nilai eigennya adalah $\lambda_1=4$ serta (dengan menggunakan
           rumus kuadrat) $\lambda_2=2+\sqrt{3}$ dan 
           $\lambda_3=2-\sqrt{3}$.
           Untuk mencari vektor eigen, tinjau sistem ini.
           \begin{equation*}
             \begin{linsys}{3}
               -x\cdot b_1  &+  &b_2          &   &               &=  &0  \\
                            &   &-x\cdot b_2  &+  &b_3            &=  &0  \\
               4\cdot b_1   &-  &17\cdot b_2  &+  &(8-x)\cdot b_3 &= &0 
             \end{linsys}
           \end{equation*}
           Menyubstitusikan $x=\lambda_1=4$ menghasilkan sistem 
           \begin{align*}
             \begin{linsys}{3}
               -4\cdot b_1  &+  &b_2          &   &           &= &0  \\
                            &   &-4\cdot b_2  &+  &b_3        &= &0  \\
               4\cdot b_1   &-  &17\cdot b_2  &+  &4\cdot b_3 &= &0 
             \end{linsys}                                              
             &\grstep{\rho_1+\rho_3}
             \begin{linsys}{3}
               -4\cdot b_1  &+  &b_2          &   &           &= &0  \\
                            &   &-4\cdot b_2  &+  &b_3        &= &0  \\
                            &   &-16\cdot b_2  &+ &4\cdot b_3 &= &0 
             \end{linsys}                                                \\
             &\grstep{-4\rho_2+\rho_3}
             \begin{linsys}{3}
               -4\cdot b_1  &+  &b_2          &   &           &= &0  \\
                            &   &-4\cdot b_2  &+  &b_3        &= &0  \\
                            &   &              &  &0          &= &0 
             \end{linsys}
           \end{align*}
           yang menghasilkan ruang eigen dan vektor eigen berikut.
           \begin{equation*}
             V_4=\set{\colvec{(1/16)\cdot b_3 \\ (1/4)\cdot b_3 \\ b_3}
                    \suchthat b_3\in\C}
             \qquad
             \colvec[r]{1 \\ 4 \\ 16}
           \end{equation*}

           Menyubstitusikan $x=\lambda_2=2+\sqrt{3}$ menghasilkan sistem 
           \begin{multline*}
             \begin{linsys}{3}
               (-2-\sqrt{3})\cdot b_1  
                      &+  &b_2          
                          &   &           &= &0  \\
                      &   &(-2-\sqrt{3})\cdot b_2  
                          &+  &b_3        &= &0  \\
               4\cdot b_1   
                      &-  &17\cdot b_2  
                          &+  &(6-\sqrt{3})\cdot b_3 &= &0 
             \end{linsys}                                   \\
             \grstep{(-4/(-2-\sqrt{3}))\rho_1+\rho_3}
             \begin{linsys}{3}
               (-2-\sqrt{3})\cdot b_1  
                      &+  &b_2          
                          &   &           &= &0  \\
                      &   &(-2-\sqrt{3})\cdot b_2  
                          &+  &b_3        &= &0  \\
                      &+  &(-9-4\sqrt{3})\cdot b_2  
                          &+  &(6-\sqrt{3})\cdot b_3 &= &0 
             \end{linsys}
           \end{multline*}
           (koefisien tengah pada persamaan ketiga sama dengan
           $(-4/(-2-\sqrt{3}))-17$; samakan penyebutnya menjadi
           $-2-\sqrt{3}$, lalu rasionalkan penyebut dengan mengalikan
           pembilang dan penyebut pecahan dengan $-2+\sqrt{3}$)
           \begin{equation*}
             \grstep{((9+4\sqrt{3})/(-2-\sqrt{3}))\rho_2+\rho_3}
             \begin{linsys}{3}
               (-2-\sqrt{3})\cdot b_1  
                      &+  &b_2          
                          &   &           &= &0  \\
                      &   &(-2-\sqrt{3})\cdot b_2  
                          &+  &b_3        &= &0  \\
                      &   &                         
                          &   &0                     &= &0 
             \end{linsys}
           \end{equation*}
           yang menghasilkan ruang eigen dan vektor eigen berikut.
           \begin{equation*}
             V_{2+\sqrt{3}}
             =\set{\colvec{(1/(2+\sqrt{3})^2)\cdot b_3  \\ 
                           (1/(2+\sqrt{3}))\cdot b_3    \\ 
                           b_3}
                    \suchthat b_3\in\C}
             \qquad
             \colvec{(1/(2+\sqrt{3})^2)  \\ 
                           (1/(2+\sqrt{3}))  \\ 
                           1}
           \end{equation*}

           Terakhir, menyubstitusikan $x=\lambda_3=2-\sqrt{3}$ menghasilkan sistem 
           \begin{multline*}
             \begin{linsys}{3}
               (-2+\sqrt{3})\cdot b_1  
                      &+  &b_2          
                          &   &           &= &0  \\
                      &   &(-2+\sqrt{3})\cdot b_2  
                          &+  &b_3        &= &0  \\
               4\cdot b_1   
                      &-  &17\cdot b_2  
                          &+  &(6+\sqrt{3})\cdot b_3 &= &0 
             \end{linsys}                                       \\
             \begin{aligned}
               &\grstep{(-4/(-2+\sqrt{3}))\rho_1+\rho_3}
               \begin{linsys}{3}
                 (-2+\sqrt{3})\cdot b_1  
                        &+  &b_2          
                            &   &           &= &0  \\
                        &   &(-2+\sqrt{3})\cdot b_2  
                            &+  &b_3        &= &0  \\
                        &   &(-9+4\sqrt{3})\cdot b_2  
                            &+  &(6+\sqrt{3})\cdot b_3 &= &0 
               \end{linsys}                                       \\
               &\grstep{((9-4\sqrt{3})/(-2+\sqrt{3}))\rho_2+\rho_3}
               \begin{linsys}{3}
                 (-2+\sqrt{3})\cdot b_1  
                        &+  &b_2          
                            &   &           &= &0  \\
                        &   &(-2+\sqrt{3})\cdot b_2  
                            &+  &b_3        &= &0  \\
                        &   &                         
                            &   &0                     &= &0 
               \end{linsys}
             \end{aligned}
           \end{multline*}
           yang memberikan ruang eigen dan vektor eigen berikut.
           \begin{equation*}
             V_{2-\sqrt{3}}
             =\set{\colvec{(1/(-2+\sqrt{3})^2)\cdot b_3  \\ 
                           (1/(2-\sqrt{3}))\cdot b_3    \\ 
                           b_3}
                    \suchthat b_3\in\C}
             \qquad
             \colvec{(1/(-2+\sqrt{3})^2)  \\ 
                           (1/(2-\sqrt{3}))  \\ 
                           1}
           \end{equation*}
      \end{exparts}
    \end{answer}
  \item 
   Untuk setiap matriks, tentukan polinomial karakteristik, nilai eigen,
   dan ruang eigen yang terkait.
   Tentukan juga multiplisitas aljabar dan geometriknya.
   \begin{exparts*}
     \partsitem 
       $\begin{mat}
          13 &-4 \\
          -4 &7
        \end{mat}$
     \partsitem
        $\begin{mat}
          1 &3 &-3 \\
         -3 &7 &-3 \\
         -6 &6 &-2
        \end{mat}$
    \partsitem
        $\begin{mat}
          2 &3 &-3 \\
          0 &2 &-3 \\
          0 &0 &1
        \end{mat}$
   \end{exparts*}
   \begin{answer}
     \begin{exparts}
       \item Polinomial karakteristiknya terfaktorkan sebagai 
         $x^2-20x+75=(x-5)(x-15)$,
         sehingga nilai eigennya adalah $\lambda_1=5$ dan~$\lambda_2=15$.
         Berikut ruang eigen yang terkait.
         \begin{equation*}
           V_5=\set{k\colvec{1 \\ 2}\suchthat k\in\C}
           \quad
           V_{15}=\set{k\colvec{-2 \\ 1}\suchthat k\in\C}
         \end{equation*}
         Untuk setiap nilai eigen, multiplisitas aljabar dan geometriknya
         sama-sama~$1$.
      \item Polinomial karakteristik $x^3 - 6x^2 + 32$ terfaktorkan menjadi
         $(x+2)(x-4)^2$.
         Nilai eigennya adalah $\lambda_1=-2$ dan $\lambda_2=4$.
         Berikut ruang eigen yang terkait. 
         \begin{equation*}
           V_{-2}=\set{k\colvec{1 \\ 1 \\ 2}\suchthat k\in\C}
           \quad
           V_4=\set{k_1\colvec{1 \\ 1 \\ 0}
                     +k_2\colvec{1 \\ 0 \\ -1} \suchthat k_1,k_2\in\C}
         \end{equation*}
         Untuk $\lambda_1=-2$, multiplisitas aljabar dan geometriknya
         sama-sama~$1$. 
         Untuk $\lambda_2=4$, multiplisitas aljabar dan geometriknya
         sama-sama~$2$.
      \item Polinomial karakteristik $x^3 - 5x^2 + 8x - 4$ terfaktorkan menjadi
         $(x-1)(x-2)^2$.
         Nilai eigennya adalah $\lambda_1=1$ dan $\lambda_2=2$.
         Berikut ruang eigen yang terkait.
         \begin{equation*}
           V_1=\set{k\colvec{-6 \\ 3 \\ 1}\suchthat k\in \C}
           \quad
           V_2=\set{k\colvec{1 \\ 0 \\ 0}\suchthat k\in\C}
         \end{equation*}
         Untuk $\lambda_1=1$, multiplisitas aljabar dan geometriknya
         sama-sama~$1$.
         Untuk $\lambda_2=2$, multiplisitas aljabarnya~$2$, tetapi 
         multiplisitas geometriknya~$1$.
     \end{exparts}
   \end{answer}
   \recommended \item
     Misalkan \( \map{t}{\polyspace_2}{\polyspace_2} \) adalah peta linear berikut.
     \begin{equation*}
      a_0+a_1x+a_2x^2\mapsto
      (5a_0+6a_1+2a_2)-(a_1+8a_2)x+(a_0-2a_2)x^2
     \end{equation*}
    Tentukan nilai eigen dan vektor eigen yang terkait.
    \begin{answer}
      Terhadap basis alami $B=\sequence{1,x,x^2}$, 
      representasi matriksnya adalah berikut ini.
      \begin{equation*}
        \rep{t}{B,B}
        =
        \begin{mat}[r]
          5  &6  &2  \\
          0  &-1 &-8 \\
          1  &0  &-2 
        \end{mat}
      \end{equation*}
      Dengan demikian, persamaan karakteristiknya adalah
      \begin{equation*}
        0
        =
        \begin{vmatrix}
          5-x  &6    &2  \\
          0    &-1-x &-8 \\
          1    &0    &-2-x 
        \end{vmatrix}
        =(5-x)(-1-x)(-2-x)-48-2\cdot(-1-x)
      \end{equation*}
      yang juga dapat ditulis sebagai
      $0=-x^3+2x^2+15x-36=-1\cdot (x+4)(x-3)^2$.
      Untuk mencari vektor eigen yang terkait, tinjau sistem ini.
      \begin{equation*}
        \begin{linsys}{3}
          (5-x)\cdot b_1 &+ &6b_2      &+ &2b_3      &= &0 \\
                         &  &(-1-x)\cdot b_2 &- &8b_3      &= &0 \\
          b_1            &  &                &+ &(-2-x)\cdot b_3 &= &0 
        \end{linsys}
      \end{equation*}
      Menyubstitusikan $\lambda_1=-4$ untuk $x$ menghasilkan
      \begin{equation*}
        \begin{linsys}{3}
                    9b_1 &+ &6b_2      &+ &2b_3      &= &0 \\
                         &  &3b_2      &- &8b_3      &= &0 \\
                     b_1 &  &          &+ &2b_3     &= &0 
        \end{linsys}
        \grstep{-(1/9)\rho_1+\rho_3}\quad
        \grstep{(2/9)\rho_2+\rho_3}\quad
        \begin{linsys}{3}
                     9b_1 &+ &6b_2      &+ &2b_3      &= &0 \\
                         &  &3b_2       &- &8b_3      &= &0  
        \end{linsys}
      \end{equation*}
      Berikut ruang eigen dan sebuah vektor eigennya.
           \begin{equation*}
             V_{-4}
             =\set{\colvec{2\cdot b_3  \\ 
                           (8/3)\cdot b_3    \\ 
                           b_3}
                    \suchthat b_3\in\C}
             \qquad
             \colvec{2  \\ 
                     8/3  \\ 
                      1}
           \end{equation*}

      Serupa dengan itu, menyubstitusikan $x=\lambda_2=3$ menghasilkan
      \begin{equation*}
        \begin{linsys}{3}
                    2b_1 &+ &6\cdot b_2      &+ &2\cdot b_3      &= &0 \\
                         &  &-4\cdot b_2      &- &8\cdot b_3      &= &0 \\
                     b_1 &  &                &- & 5\cdot b_3     &= &0 
        \end{linsys}
        \grstep{-(1/2)\rho_1+\rho_3}
        \repeatedgrstep{-(3/4)\rho_2+\rho_3}
        \begin{linsys}{3}
                     2b_1 &+ &6\cdot b_2      &+ &2\cdot b_3      &= &0 \\
                         &  &-4\cdot b_2       &- &8\cdot b_3      &= &0  
        \end{linsys}
      \end{equation*}
      dengan ruang eigen dan vektor eigen berikut.
           \begin{equation*}
             V_{3}
             =\set{\colvec{5\cdot b_3  \\ 
                           -2\cdot b_3    \\ 
                           b_3}
                    \suchthat b_3\in\C}
             \qquad
             \colvec[r]{5  \\ 
                     -2  \\ 
                       1}
           \end{equation*}
    \end{answer}
   \item 
     Tentukan nilai eigen dan vektor eigen dari
     pemetaan \( \map{t}{\matspace_2}{\matspace_2} \) berikut.
     \begin{equation*}
         \begin{mat}
            a  &b  \\
            c  &d
         \end{mat}
       \mapsto
         \begin{mat}
           2c    &a+c  \\
           b-2c  &d
         \end{mat}
     \end{equation*}
     \begin{answer}
        $\lambda=1,
          \begin{mat}
                   0  &0  \\
                   0  &1
          \end{mat} \text{ dan }
          \begin{mat}[r]
                   2  &3  \\
                   1  &0
          \end{mat}$,           
          $\lambda=-2,
          \begin{mat}[r]
                  -1  &0  \\
                   1  &0
          \end{mat}$,
          $\lambda=-1,
          \begin{mat}[r]
                  -2  &1  \\
                   1  &0
          \end{mat}$  
       \end{answer}
   \recommended \item 
     Tentukan nilai eigen dan vektor eigen yang terkait dari
     operator diferensiasi
     \( \map{d/dx}{\polyspace_3}{\polyspace_3} \).
     \begin{answer}
       Tetapkan basis alami $B=\sequence{1,x,x^2,x^3}$. 
       Aksi pemetaan ini adalah $1\mapsto 0$, $x\mapsto 1$, $x^2\mapsto 2x$,
       dan $x^3\mapsto 3x^2$, sehingga representasinya mudah dihitung.
       \begin{equation*}
         T=\rep{d/dx}{B,B}=
         \begin{mat}[r]
           0  &1  &0  &0  \\
           0  &0  &2  &0  \\
           0  &0  &0  &3  \\
           0  &0  &0  &0
         \end{mat}_{B,B}
       \end{equation*}
       Kita menentukan nilai eigennya dengan perhitungan berikut.
       \begin{equation*}
         0=\deter{T-xI}=
         \begin{vmatrix}
           -x &1  &0  &0  \\
           0  &-x &2  &0  \\
           0  &0  &-x &3  \\
           0  &0  &0  &-x          
         \end{vmatrix}
         =x^4
       \end{equation*}
       Jadi, pemetaan ini hanya memiliki nilai eigen $\lambda=0$.
       Untuk menentukan vektor eigen yang terkait, kita selesaikan
       \begin{equation*}
         \begin{mat}[r]
           0  &1  &0  &0  \\
           0  &0  &2  &0  \\
           0  &0  &0  &3  \\
           0  &0  &0  &0
         \end{mat}_{B,B}
         \colvec{b_1 \\ b_2 \\ b_3 \\ b_4}_B
         =0\cdot\colvec{b_1 \\ b_2 \\ b_3 \\ b_4}_B
         \qquad\Longrightarrow\qquad
         \text{$b_2=0$, $b_3=0$, $b_4=0$}          
       \end{equation*}
       sehingga diperoleh ruang eigen berikut.
       \begin{equation*}
         \set{\colvec{b_1 \\ 0 \\ 0 \\ 0}_B
               \suchthat b_1\in\C}
         =\set{b_1+0\cdot x+0\cdot x^2+0\cdot x^3
               \suchthat b_1\in\C}
         =\set{b_1
               \suchthat b_1\in\C}
       \end{equation*}
      \end{answer}
   \item Buktikan bahwa 
     nilai eigen suatu matriks segitiga (atas atau bawah)
     adalah entri-entri pada diagonalnya.
     \begin{answer}
       Determinan matriks segitiga $T-xI$ adalah hasil kali 
       entri-entri pada diagonalnya, sehingga dapat difaktorkan menjadi hasil kali 
       suku-suku $t_{i,i}-x$.
     \end{answer}
  \recommended \item Matriks berikut memiliki nilai-nilai eigen yang
    berbeda.
    \begin{equation*}
     \begin{mat}
      1  &2  &1 \\
      6  &-1 &0 \\
      -1 &-2 &-1
     \end{mat}
    \end{equation*}
    \begin{exparts}
      \partsitem Diagonalkan matriks ini.
      \partsitem Tentukan basis yang terhadapnya matriks ini memiliki
        representasi diagonal tersebut.
      \partsitem Gambarkan diagramnya.
         Tentukan matriks $P$ dan $P^{-1}$ yang menghasilkan perubahan basis itu. 
    \end{exparts}
    \begin{answer}
     \begin{exparts}
      \partsitem 
      Perhitungan nilai eigen memberikan $\lambda_1=3$, $\lambda_2=0$,
      dan $\lambda_3=-4$.
      Jadi, diagonalisasi matriks tersebut adalah sebagai berikut.
      \begin{equation*}
        \begin{mat}
          3 &0 &0 \\
          0 &0 &0 \\
          0 &0 &-4
        \end{mat}
      \end{equation*}
     \partsitem
      Tentukan satu vektor eigen yang terkait dengan setiap nilai eigen, lalu
      gunakan vektor-vektor itu untuk membentuk suatu basis.

      Untuk $\lambda_1=3$, kita selesaikan sistem berikut.
      \begin{equation*}
        \begin{mat}
          1-3  &2    &1 \\
          6    &-1-3 &0 \\
         -1    &-2   &-1-3
        \end{mat}
        \colvec{x \\ y \\ z}
        =\colvec{0 \\ 0 \\ 0}
      \end{equation*}
      Metode Gauss sudah memadai.
      % sage: M = matrix(QQ, [[-2,2,1,0], [6,-4,0,0], [-1,-2,-4,0]])
      % sage: gauss_method(M)
      % [-2  2  1  0]
      % [ 6 -4  0  0]
      % [-1 -2 -4  0]
      % take 3 times row 1 plus row 2
      % take -1/2 times row 1 plus row 3
      % [  -2    2    1    0]
      % [   0    2    3    0]
      % [   0   -3 -9/2    0]
      % take 3/2 times row 2 plus row 3
      % [-2  2  1  0]
      % [ 0  2  3  0]
      % [ 0  0  0  0]
      \begin{equation*}
        \begin{amat}{3}
          -2 &2  &1  &0  \\
           6 &-4 &0  &0  \\
          -1 &-2 &-4 &0
        \end{amat}
        \grstep[-(1/2)\rho_1+\rho_3]{3\rho_1+\rho_2}
        \grstep{(3/2)\rho_2+\rho_3}
        \begin{amat}{3}
          -2 &2 &1   &0  \\
           0 &2 &3   &0  \\
           0 &0 &0   &0
        \end{amat}
      \end{equation*}
      Parametrisasi memberikan ruang solusi berikut.
      \begin{equation*}
        S_1=\set{\colvec{-1 \\ -3/2 \\ 1}\cdot z \suchthat z\in\C}
      \end{equation*}

      Berikut adalah sistem untuk $\lambda_2=0$.
      \begin{equation*}
        \begin{mat}
          1    &2    &1 \\
          6    &-1 &0 \\
         -1    &-2   &-1
        \end{mat}
        \colvec{x \\ y \\ z}
        =\colvec{0 \\ 0 \\ 0}
      \end{equation*}
      Metode Gauss 
      % sage: M = matrix(QQ, [[1,2,1,0], [6,-1,0,0], [-1,-2,-1,0]])
      % sage: gauss_method(M)
      % [ 1  2  1  0]
      % [ 6 -1  0  0]
      % [-1 -2 -1  0]
      % take -6 times row 1 plus row 2
      % take 1 times row 1 plus row 3
      % [  1   2   1   0]
      % [  0 -13  -6   0]
      % [  0   0   0   0]
      \begin{equation*}
        \begin{amat}{3}
           1 &2  &1  &0  \\
           6 &-1 &0  &0  \\
          -1 &-2 &-1 &0
        \end{amat}
        \grstep[\rho_1+\rho_3]{-6\rho_1+\rho_2}
        \begin{amat}{3}
           1 &2 &1   &0  \\
           0 &-13 &-6   &0  \\
           0 &0 &0   &0
        \end{amat}
      \end{equation*}
      yang diikuti dengan parametrisasi memberikan ruang solusi berikut.
      \begin{equation*}
        S_2=\set{\colvec{-1/13 \\ -6/13 \\ 1}\cdot z \suchthat z\in\C}
      \end{equation*}

      Sistem untuk $\lambda_3=-4$ dikerjakan dengan cara yang sama.
      \begin{equation*}
        \begin{mat}
          1-(-4)    &2    &1 \\
          6    &-1-(-4) &0 \\
         -1    &-2   &-1-(-4)
        \end{mat}
        \colvec{x \\ y \\ z}
        =\colvec{0 \\ 0 \\ 0}
      \end{equation*}
      Metode Gauss 
      % sage: M = matrix(QQ, [[5,2,1,0], [6,3,0,0], [-1,-2,3,0]])
      % sage: gauss_method(M)
      % [ 5  2  1  0]
      % [ 6  3  0  0]
      % [-1 -2  3  0]
      % take -6/5 times row 1 plus row 2
      % take 1/5 times row 1 plus row 3
      % [   5    2    1    0]
      % [   0  3/5 -6/5    0]
      % [   0 -8/5 16/5    0]
      % take 8/3 times row 2 plus row 3
      % [   5    2    1    0]
      % [   0  3/5 -6/5    0]
      % [   0    0    0    0]
      \begin{equation*}
        \begin{amat}{3}
           5 &2  &1  &0  \\
           6 &3 &0  &0  \\
          -1 &-2 &3 &0
        \end{amat}
        \grstep[(1/5)\rho_1+\rho_3]{-(6/5)\rho_1+\rho_2}
        \grstep{(8/3)\rho_2+\rho_3}
        \begin{amat}{3}
           5 &2 &1   &0  \\
           0 &3/5 &-6/5   &0  \\
           0 &0 &0   &0
        \end{amat}
      \end{equation*}
      dan parametrisasi memberikan hasil berikut.
      \begin{equation*}
        S_3=\set{\colvec{-1 \\ 2 \\ 1}\cdot z \suchthat z\in\C}
      \end{equation*}

      Jadi, matriks tersebut memiliki bentuk diagonal itu ketika merepresentasikan
      transformasi yang sama terhadap pasangan basis $D,D$ berikut.
      \begin{equation*}
        D=\sequence{\colvec{-1 \\ -3/2 \\ 1},
                    \colvec{-1/13 \\ -6/13 \\ 1},
                    \colvec{-1 \\ 2 \\ 1}}
      \end{equation*}
     \partsitem
     % sage: M = matrix(QQ, [[1,2,1], [6,-1,0], [-1,-2,-1]])
     % sage: M.eigenvalues()
     % [3, 0, -4]
     % sage: P = matrix(QQ, [[-1,-1/13,-1], [-3/2,-6/13,2], [1,1,1]])
     % sage: P.inverse()*M*P
     % [ 3  0  0]
     % [ 0  0  0]
     % [ 0  0 -4]
     % sage: P.inverse()
     % [-16/21   -2/7  -4/21]
     % [ 13/12      0  13/12]
     % [ -9/28    2/7   3/28]
       Namakan matriks tersebut~$T$.
       Dengan mengambil ruang domain dan kodomain sebagai $\Re^3$ serta
       $T=\rep{t}{\stdbasis_3,\stdbasis_3}$, kita memperoleh diagram berikut.
       \begin{equation*}
       \begin{CD}
          \C^3_{\wrt{\stdbasis_3}}            @>t>T>      \C^3_{\wrt{\stdbasis_3}}       \\
          @V{\scriptstyle\identity} VV           @V{\scriptstyle\identity} VV \\
          \C^3_{\wrt{D}}             @>t>D>  \C^3_{\wrt{D}}
        \end{CD}
        \end{equation*}
        Dari diagram tersebut kita memperoleh $D=P^{-1}TP$, dengan
        $P=\rep{\identity}{D,\stdbasis_3}$.
        Merepresentasikan setiap vektor dalam $D$ terhadap basis standar
        mudah dilakukan karena, terhadap basis standar, setiap vektor
        merepresentasikan dirinya sendiri.
        \begin{multline*}
          \rep{\colvec{-1 \\ -3/2 \\ 1}}{\stdbasis_3}=\colvec{-1 \\ -3/2 \\ 1}
          \quad
          \rep{\colvec{-1/13 \\ -6/13 \\ 1}}{\stdbasis_3}=\colvec{-1/13 \\ -6/13 \\ 1}
          \\
          \rep{\colvec{-1 \\ 2 \\ 1}}{\stdbasis_3}=\colvec{-1 \\ 2 \\ 1}
        \end{multline*}
        Jadi, matriks itu cukup menyandingkan ketiga vektor basis tersebut.
        \begin{equation*}
          P=
          \begin{mat}
            -1   &-1/13 &-1 \\
            -3/2 &-6/13 &2  \\
            1    &1     &1
          \end{mat}
         \quad\text{sehingga}\quad
         P^{-1}=
         \begin{mat}
           -16/21 &-2/7 &-4/21 \\
           13/12  &0    &13/12 \\
           -9/28  &2/7  &3/28
         \end{mat}
        \end{equation*}





     \end{exparts}
    \end{answer}
  \recommended \item 
    Tentukan rumus polinomial karakteristik suatu matriks $\nbyn{2}$.
    \begin{answer}
      Cukup ekspansikan determinan $T-xI$.
      \begin{equation*}
        \begin{vmatrix}
          a-x  &c  \\
          b    &d-x
        \end{vmatrix}
        =(a-x)(d-x)-bc
        =x^2+(-a-d)\cdot x +(ad-bc)
      \end{equation*}
    \end{answer}
  \item \label{exer:CharPolyTransWellDefed}
    Buktikan bahwa 
    polinomial karakteristik suatu transformasi terdefinisi dengan baik.
    \begin{answer}
      Sebarang dua representasi transformasi tersebut saling serupa, dan
      matriks-matriks yang serupa memiliki polinomial karakteristik yang sama.  
    \end{answer}
  \item Buktikan atau sangkal: jika semua nilai eigen suatu matriks adalah $0$, 
    maka matriks itu harus merupakan matriks nol.
    \begin{answer}
      Pernyataan itu tidak benar.
      Semua nilai eigen matriks berikut adalah $0$.
      \begin{equation*}
        \begin{mat}[r]
          0  &1  \\
          0  &0
        \end{mat}
      \end{equation*}
    \end{answer}
  \recommended \item 
    \begin{exparts}
      \partsitem Tunjukkan bahwa sebarang vektor tak nol dalam sebarang 
        ruang vektor tak trivial dapat menjadi vektor eigen.
        Artinya, untuk \( \vec{v}\neq\zero \) dalam ruang tak trivial \( V \),
        tunjukkan bahwa terdapat transformasi \( \map{t}{V}{V} \) dengan
        nilai eigen skalar \( \lambda\in\Re \) sedemikian sehingga
        \( \vec{v}\in V_{\lambda} \).
      \partsitem Bagaimana jika suatu skalar \( \lambda \) juga diberikan?
        Dapatkah sebarang anggota tak nol dari sebarang ruang vektor tak trivial
        menjadi vektor eigen yang terkait dengan \( \lambda \)?
    \end{exparts}
    \begin{answer}
      \begin{exparts}
        \partsitem Gunakan \( \lambda=1 \) dan pemetaan identitas.
        \partsitem Ya. Gunakan transformasi yang mengalikan semua 
          vektor dengan skalar \( \lambda \).
      \end{exparts}  
     \end{answer}
  \recommended \item 
    Misalkan \( \map{t}{V}{V} \) dan \( T=\rep{t}{B,B} \).
    Buktikan bahwa vektor-vektor eigen \( T \) yang terkait dengan \( \lambda \)
    adalah vektor-vektor tak nol dalam kernel pemetaan yang direpresentasikan
    (terhadap pasangan basis yang sama) oleh \( T-\lambda I \).
    \begin{answer}
      Jika $t(\vec{v})=\lambda\cdot\vec{v}$, maka 
      $\vec{v}\mapsto\zero$ di bawah pemetaan $t-\lambda\cdot\identity$.
    \end{answer}
  \item 
    Buktikan bahwa jika $a,\ldots,\,d$ semuanya bilangan bulat dan
    \( a+b=c+d \), maka
    \begin{equation*}
      \begin{mat}
         a  &b  \\
         c  &d
      \end{mat}
    \end{equation*}
    memiliki nilai eigen bulat, yaitu \( a+b \) dan \( a-c \).
    \begin{answer}
      Persamaan karakteristiknya 
      \begin{equation*}
        0=
        \begin{vmatrix}
          a-x  &b  \\
          c    &d-x 
        \end{vmatrix}
        =(a-x)(d-x)-bc
      \end{equation*}
      menyederhana menjadi $x^2+(-a-d)\cdot x + (ad-bc)$.
      Pemeriksaan bahwa nilai $x=a+b$ dan $x=a-c$ memenuhi persamaan tersebut 
      (dengan syarat $a+b=c+d$) bersifat rutin. 
    \end{answer}
  \recommended \item
    Buktikan bahwa jika \( T \) nonsingular dan memiliki nilai eigen
    \( \lambda_1,\dots,\lambda_n \), maka \( T^{-1} \) memiliki nilai eigen
    \( 1/\lambda_1,\dots,1/\lambda_n \).
    Apakah konversnya benar?
    \begin{answer}
      Perhatikan ruang eigen $V_{\lambda}$.
      Sebarang $\vec{w}\in V_{\lambda}$ merupakan bayangan
      $\vec{w}=\lambda\cdot\vec{v}$ dari suatu $\vec{v}\in V_{\lambda}$ 
      (yakni, $\vec{v}=(1/\lambda)\cdot\vec{w}$).
      Jadi, pada $V_{\lambda}$ (yang merupakan subruang tak trivial), 
      aksi $t^{-1}$ adalah 
      $t^{-1}(\vec{w})=\vec{v}=(1/\lambda)\cdot\vec{w}$,
      sehingga $1/\lambda$ merupakan nilai eigen $t^{-1}$.
      Konversnya juga benar, dengan menerapkan hasil ini pada $T^{-1}$ dan
      menggunakan $(T^{-1})^{-1}=T$.
    \end{answer}
  \recommended \item
    Misalkan \( T \) berukuran \( \nbyn{n} \) dan \( c,d \) adalah skalar.
    \begin{exparts}
      \partsitem Buktikan bahwa jika \( T \) memiliki nilai eigen 
        \( \lambda \) dengan vektor eigen terkait \( \vec{v} \), maka
        \( \vec{v} \) merupakan vektor eigen dari \( cT+dI \) yang terkait
        dengan nilai eigen \( c\lambda+d \).
      \partsitem Buktikan bahwa jika \( T \) dapat didiagonalkan, maka
        \( cT+dI \) juga dapat didiagonalkan.
    \end{exparts}
    \begin{answer}
      \begin{exparts}
        \partsitem Kita mempunyai 
          $(cT+dI)\vec{v}=cT\vec{v}+dI\vec{v}=c\lambda\vec{v}+d\vec{v}
               =(c\lambda+d)\cdot \vec{v}$.
        \partsitem Misalkan $S=PTP^{-1}$ diagonal.
          Maka $P(cT+dI)P^{-1}=P(cT)P^{-1}+P(dI)P^{-1}
                 =cPTP^{-1}+dI=cS+dI$ juga diagonal.
      \end{exparts}
    \end{answer}
  \recommended \item
    Tunjukkan bahwa \( \lambda \) merupakan nilai eigen \( T \) jika dan hanya
    jika pemetaan yang direpresentasikan oleh \( T-\lambda I \) bukan isomorfisme.
    \begin{answer}
      Skalar $\lambda$ merupakan nilai eigen jika dan hanya jika transformasi
      $t-\lambda \identity$ singular.
      Suatu transformasi singular jika dan hanya jika transformasi itu bukan
      isomorfisme (yakni, suatu transformasi merupakan isomorfisme jika dan
      hanya jika transformasi itu nonsingular).
    \end{answer}
  \item \cite{Strang} 
    \begin{exparts}
      \partsitem Tunjukkan bahwa jika \( \lambda \) merupakan nilai eigen
         \( A \), maka \( \lambda^k \) merupakan nilai eigen \( A^k \)
         untuk setiap bilangan bulat positif \( k \).
      \partsitem Apa yang salah dengan bukti generalisasi berikut?
         ``Jika \( \lambda \) merupakan nilai eigen \( A \) dan \( \mu \)
         merupakan nilai eigen \( B \), maka \( \lambda\mu \) merupakan
         nilai eigen \( AB \), sebab jika \( A\vec{x}=\lambda\vec{x} \) dan
         \( B\vec{x}=\mu\vec{x} \), maka
         \( AB\vec{x}=A\mu\vec{x}=\mu A\vec{x}=\mu\lambda\vec{x} \).''
    \end{exparts}
    \begin{answer}
      \begin{exparts}
        \partsitem Jika nilai eigen $\lambda$ terkait dengan vektor eigen
          $\vec{x}$, maka
          $A^k\vec{x}=A\cdots A\vec{x}=A^{k-1}\lambda\vec{x}
            =\lambda A^{k-1}\vec{x}=\cdots=\lambda^k\vec{x}$.
          (Rincian lengkapnya memerlukan induksi pada $k$.)
        \partsitem Vektor eigen yang terkait dengan $\lambda$
          belum tentu merupakan vektor eigen yang terkait dengan $\mu$.
      \end{exparts}
    \end{answer}
  \item 
    Apakah matriks-matriks yang ekuivalen memiliki nilai eigen yang sama?
    \begin{answer}
      Tidak.
      Berikut adalah dua matriks berukuran sama dan berperingkat sama,
      tetapi memiliki nilai eigen yang berbeda.
      \begin{equation*}
        \begin{mat}[r]
          1  &0  \\
          0  &1
        \end{mat}
        \qquad
        \begin{mat}[r]
          1  &0  \\
          0  &2
        \end{mat}
      \end{equation*}
    \end{answer}
  \item 
    Tunjukkan bahwa matriks persegi berentri riil dengan jumlah baris ganjil
    memiliki sekurang-kurangnya satu nilai eigen riil.
    \begin{answer}
      Polinomial karakteristiknya berderajat ganjil sehingga 
      memiliki sekurang-kurangnya satu akar riil.  
    \end{answer}
  \item 
    Diagonalkan.
    \begin{equation*}
       \begin{mat}[r]
         -1  &2  &2  \\
          2  &2  &2  \\
         -3  &-6 &-6
       \end{mat}
    \end{equation*}
    \begin{answer}
      Polinomial karakteristik $x^3+5x^2+6x$ memiliki akar-akar berbeda
      \( \lambda_1=0 \), \( \lambda_2=-2 \), dan \( \lambda_3=-3 \).
      Jadi, matriks tersebut dapat didiagonalkan menjadi bentuk berikut.
      \begin{equation*}
         \begin{mat}[r]
            0  &0  &0  \\
            0  &-2 &0  \\
            0  &0  &-3
         \end{mat}
      \end{equation*}    
    \end{answer}
  \item 
    Misalkan \( P \) merupakan matriks nonsingular berukuran \( \nbyn{n} \).
    Tunjukkan bahwa pemetaan
    \definend{transformasi keserupaan}\index{transformasi keserupaan}
    \( \map{t_P}{\matspace_{\nbyn{n}}}{\matspace_{\nbyn{n}}} \)
    yang mengirimkan \( T\mapsto PTP^{-1} \)
    merupakan isomorfisme.
    \begin{answer}
      Kita harus menunjukkan bahwa pemetaan itu satu-satu dan pada, serta
      mempertahankan operasi penjumlahan matriks dan perkalian skalar.

      Untuk menunjukkan bahwa pemetaan itu satu-satu, misalkan
      $t_P(T)=t_P(S)$, yakni $PTP^{-1}=PSP^{-1}$. Mengalikan kedua ruas
      dari kiri dengan $P^{-1}$ dan dari kanan dengan $P$ memberikan $T=S$.
      Untuk menunjukkan bahwa pemetaan itu pada, ambil
      $S\in\matspace_{\nbyn{n}}$ dan perhatikan bahwa
      $S=t_P(P^{-1}SP)$.

      Pemetaan $t_P$ mempertahankan penjumlahan matriks karena
      $t_P(T+S)=P(T+S)P^{-1}=(PT+PS)P^{-1}=PTP^{-1}+PSP^{-1}
      =t_P(T)+t_P(S)$, berdasarkan sifat perkalian dan penjumlahan matriks
      yang telah kita pelajari.
      Perkalian skalar serupa:
      $t_P(cT)=P(c\cdot T)P^{-1}=c\cdot (PTP^{-1})=c\cdot t_P(T)$.
    \end{answer}
  \puzzle \item 
    \cite{MathMag67p232}
    Tunjukkan bahwa jika \( A \) merupakan matriks persegi berordo \( n \)
    dan jumlah entri pada setiap baris (kolom) adalah \( c \), maka \( c \)
    merupakan akar karakteristik \( A \).
    (``Akar karakteristik'' adalah sinonim untuk nilai eigen.)\index{karakteristik!akar}\index{akar!karakteristik}
    \begin{answer}
      \answerasgiven %
      Jika variabel dalam polinomial karakteristik \( A \) ditetapkan sama
      dengan \( c \), penjumlahan \( (n-1) \) baris (kolom) pertama ke
      baris (kolom) ke-\( n \) menghasilkan determinan yang baris (kolom)
      ke-\( n \)-nya nol.
      Jadi, \( c \) merupakan akar karakteristik \( A \).  
    \end{answer}
\index{keserupaan|)}
\end{exercises}
